\documentclass[a4paper]{article} %tạo một bản báo cáo
\usepackage[utf8]{inputenc}
\usepackage[T5]{fontenc} %để sử dụng tiếng việt
\usepackage[fontsize=11pt]{scrextend} %set font=13pt
\linespread{1.2}
\usepackage[paperheight=29.7cm,paperwidth=21cm,right=2cm,left=1.8cm,top=2cm,bottom=2.5cm]{geometry}
\usepackage{amsmath}
\usepackage{amsmath, amssymb, txfonts}
\usepackage{listings}
\usepackage{tabularx}
\usepackage{array}
\usepackage{diagbox}
\usepackage{minted}
\usepackage{vhistory}
\usepackage{indentfirst}
\usepackage{multirow}
\usepackage{graphicx} %thư viện chèn ảnh
\usepackage{float} % set vị trí chèn ảnh
\usepackage{tikz} % thư viện tạo khung bìa
\usetikzlibrary{calc} % thư viện tikz
\renewcommand{\figurename}{\fontsize{12pt}{0pt}\selectfont \bfseries Hình}
\renewcommand{\listfigurename}{Danh mục hình ảnh}
\renewcommand{\listtablename}{Danh mục bảng}
\renewcommand{\thefigure}{\thesection.\arabic{figure}}
\usepackage{caption}
\usepackage{enumerate} % để tùy chỉnh itemize
\usepackage{array}    % để dùng m{width}
\captionsetup[figure]{labelsep=space}
\newtheorem{theorem}{Định lý}[section]
\newtheorem{defn}{Định nghĩa}[section]
\newtheorem{vd}{Ví dụ}[section]
\newtheorem{proof}{Chứng minh}
\usepackage{adjustbox}
\usepackage{esint}
\usepackage{fancyhdr}
\usepackage[colorlinks=true, linkcolor=black, citecolor=black, urlcolor=blue]{hyperref}
\usepackage{placeins}
\usepackage{subcaption}
\usepackage{microtype}
\usepackage{booktabs}
\usepackage{longtable}
\usepackage{lastpage}
\usepackage{enumerate}
\usepackage{tcolorbox}
\usepackage{changepage}
\usepackage{hyperref}
\usepackage[vietnamese]{babel}
\usepackage[
backend=biber,
style=numeric,
sorting=ynt
]{biblatex}
\DeclareLanguageMapping{vietnamese}{english}

% ================= THÔNG TIN CHUNG CỦA BÁO CÁO =================
\newcommand{\CourseName}{Hệ quản trị cơ sở dữ liệu}
\newcommand{\CourseNameUpper}{HỆ QUẢN TRỊ CƠ SỞ DỮ LIỆU}
\newcommand{\CourseCode}{CO3021}
\newcommand{\ReportTitle}{XÂY DỰNG ỨNG DỤNG ĐỌC VÀ SÁNG TÁC TRUYỆN TRỰC TUYẾN NOVELLE TRÊN NỀN MS SQL SERVER VÀ REDIS}
\newcommand{\Lecturer}{ThS. Lê Thị Bảo Thu} % TODO: điền tên GVHD
\newcommand{\ClassCode}{L01} % TODO: kiểm tra lại lớp

% ================= THÀNH VIÊN NHÓM =================
\newcommand{\StudentOneName}{Huỳnh Minh Tiến}
\newcommand{\StudentOneID}{2313425}
\newcommand{\StudentOneEmail}{tien.huynhminhcse@hcmut.edu.vn}
\newcommand{\StudentTwoName}{Nguyễn Khánh Trinh}
\newcommand{\StudentTwoID}{2313579}
\newcommand{\StudentTwoEmail}{trinh.nguyenkhanh@hcmut.edu.vn}
\newcommand{\StudentThreeName}{Đặng Quang Phong}
\newcommand{\StudentThreeID}{2312613}
\newcommand{\StudentThreeEmail}{phong.danghandsomek@hcmut.edu.vn}
\newcommand{\StudentFourName}{Nguyễn Hoàng Minh Quân}
\newcommand{\StudentFourID}{2312836}
\newcommand{\StudentFourEmail}{quan.nguyen090505@hcmut.edu.vn}
\newcommand{\StudentFiveName}{Trần Nguyễn Phú Nghĩa}
\newcommand{\StudentFiveID}{2312284}
\newcommand{\StudentFiveEmail}{nghia.trannguyenphu@hcmut.edu.vn}

% ================= ĐỊNH DẠNG TIÊU ĐỀ CHƯƠNG =================
% Mỗi \section là một chương: "CHƯƠNG 1. GIỚI THIỆU ĐỀ TÀI", bắt đầu ở trang mới
\usepackage{titlesec}
\titleformat{\section}{\normalfont\Large\bfseries}{CHƯƠNG~\thesection.}{0.5em}{\MakeUppercase}
\newcommand{\sectionbreak}{\clearpage}

% \renewcommand\thesection{\Roman{section}.}            % 1, 2, 3
% \renewcommand\thesubsection{\arabic{subsection}/}        % a, b, c
% \renewcommand\thesubsubsection{\alph{subsubsection}/}% 1, 2, 3 (lại số nhỏ)

% --- Nhưng vẫn muốn figure/table/equation dùng số thường ---
\makeatletter
\renewcommand{\thefigure}{\arabic{section}.\arabic{subsection}.\arabic{figure}}
\renewcommand{\thetable}{\arabic{section}.\arabic{subsection}.\arabic{table}}
\renewcommand{\theequation}{\arabic{section}.\arabic{subsection}.\arabic{equation}}
% \geometry{left=3cm,right=2.5cm,top=2.5cm,bottom=2.5cm}
\captionsetup{font=small,labelfont=bf}

% Reset lại bộ đếm hình/bảng/phương trình mỗi khi qua section mới
\@addtoreset{figure}{subsection}
\@addtoreset{table}{subsection}
\@addtoreset{equation}{subsection}
\makeatother
\addbibresource{references.bib}

\setlength{\headheight}{40pt}
\pagestyle{fancy}
\fancyhead{} % clear all header fields
\fancyhead[L]{
 \begin{tabular}{rl}
    \begin{picture}(25,15)(0,0)
    \put(0,-8){\includegraphics[width=8mm, height=8mm]{image/HCMUT_official_logo.png}}
    %\put(0,-8){\epsfig{width=10mm,figure=hcmut.eps}}
   \end{picture}&
	%\includegraphics[width=8mm, height=8mm]{01_logobachkhoa.png} & %
	\begin{tabular}{l}
		\textbf{\bf \ttfamily ĐẠI HỌC QUỐC GIA THÀNH PHỐ HỒ CHÍ MINH}\\
		\textbf{\bf \ttfamily TRƯỜNG ĐẠI HỌC BÁCH KHOA}
	\end{tabular} 	
 \end{tabular}
}
\fancyhead[R]{
	\begin{tabular}{l}
		\tiny \bf \\
		\tiny \bf 
	\end{tabular}  }
\fancyfoot{} % clear all footer fields
\fancyfoot[L]{\scriptsize \ttfamily Báo cáo Bài tập lớn - Môn: \CourseName{} (\CourseCode)}
\fancyfoot[R]{\scriptsize \ttfamily Trang {\thepage}/\pageref{LastPage}}
\renewcommand{\headrulewidth}{0.3pt}
\renewcommand{\footrulewidth}{0.3pt}



\begin{document}
    \begin{titlepage}
        \begin{tikzpicture}[overlay,remember picture]
			\draw [line width=3pt]
			($ (current page.north west) + (2.0cm,-2.0cm) $)
			rectangle
			($ (current page.south east) + (-2.0cm,2.5cm) $);
			\draw [line width=0.5pt]
			($ (current page.north west) + (2.1cm,-2.1cm) $)
			rectangle
			($ (current page.south east) + (-2.1cm,2.6cm) $);
		\end{tikzpicture}
		\begin{center}
			\vspace{-6pt} \fontsize{13pt}{0pt}\selectfont ĐẠI HỌC QUỐC GIA THÀNH PHỐ HỒ CHÍ MINH \\
			\textbf{\fontsize{16pt}{0pt}\selectfont TRƯỜNG ĐẠI HỌC BÁCH KHOA}\\
            \textbf{\fontsize{16pt}{0pt}\selectfont KHOA KHOA HỌC VÀ KỸ THUẬT MÁY TÍNH}
			\vspace{1cm}
			\begin{figure}[H]
				\centering
				\includegraphics[width=4.318cm,height=4.318cm]{image/HCMUT_official_logo.png}
			\end{figure}
			\vspace{0.5cm}
            \rule{0.7\linewidth}{1pt} \\[0.3cm] % Đường thẳng trên
			\textbf{\fontsize{16pt}{0pt}\selectfont BÁO CÁO BÀI TẬP LỚN}\\
			\textbf{\fontsize{16pt}{0pt}\selectfont MÔN: \CourseNameUpper{} (\CourseCode)}\\
                \vspace{0.5cm}
                \parbox{0.9\linewidth}{\centering\textbf{\fontsize{16pt}{20pt}\selectfont \ReportTitle}}\\
            \vspace{0.5cm}
            \rule{0.7\linewidth}{1pt} \\[1cm] % Đường thẳng dưới
		\end{center}

    \begin{table}[h]
    \centering
        \begin{tabular}{rl}
        \hspace{3 cm}\textbf{GVHD}:
        & \Lecturer\\
        \textbf{Lớp}: & \ClassCode\\
        \textbf{Sinh viên thực hiện}: & \StudentOneName{} - \StudentOneID\\
        & \StudentTwoName{} - \StudentTwoID\\
        & \StudentThreeName{} - \StudentThreeID\\
        & \StudentFourName{} - \StudentFourID\\
        & \StudentFiveName{} - \StudentFiveID\\
        \end{tabular}
    \end{table}
        \begin{center}
			\fontsize{12pt}{0pt}\selectfont\ Thành phố Hồ Chí Minh, tháng \number\month\space năm \number\year\\
		\end{center}
    \end{titlepage}
    \cleardoublepage

    \begin{center}
        \fontsize{13pt}{0pt}\selectfont \textbf{DANH SÁCH THÀNH VIÊN} \\
        
    \end{center}
    \hspace{-1cm}

    \begin{table}[H]
        \centering
        \small
        \renewcommand{\arraystretch}{1.4}
        \setlength{\tabcolsep}{8pt}
        \begin{tabular}{|c|m{4.2cm}|c|m{6.2cm}|>{\centering\arraybackslash}m{1.5cm}|}
            \hline
            \textbf{STT} & \textbf{Họ và tên} & \textbf{MSSV} & \textbf{Nhiệm vụ} & \textbf{\% Hoàn thành}\\
            \hline
            1 & \StudentOneName & \StudentOneID & & 100 \\
            \hline
            2 & \StudentTwoName & \StudentTwoID & & 100 \\
            \hline
            3 & \StudentThreeName & \StudentThreeID & & 100 \\
            \hline
            4 & \StudentFourName & \StudentFourID & & 100 \\
            \hline
            5 & \StudentFiveName & \StudentFiveID & & 100 \\
            \hline
        \end{tabular}
    \end{table}
    \textbf{Ghi chú:} 


    \cleardoublepage
    \tableofcontents
    \cleardoublepage
    \phantomsection
    \addcontentsline{toc}{section}{Danh mục bảng}
    \listoftables
    \cleardoublepage
    \phantomsection
    \addcontentsline{toc}{section}{Danh mục hình ảnh}
    \listoffigures
    \cleardoublepage
    \begin{versionhistory}
        \vhEntry{1.0}{2026.11.15}{TienHM}{Khởi tạo cấu trúc báo cáo môn \CourseName{}}
    \end{versionhistory}
    \cleardoublepage

    \section{Giới thiệu đề tài}
\label{sec:introduction}

Ngày nay, gần như mọi ứng dụng mà chúng ta sử dụng hằng ngày, từ ứng dụng ngân hàng, trang thương mại điện tử cho đến mạng xã hội, đều được xây dựng xoay quanh dữ liệu. Mỗi lần người dùng đăng nhập, tìm kiếm một sản phẩm hay đọc một bài viết, phía sau đều có hàng loạt thao tác lưu trữ và truy xuất dữ liệu diễn ra. Để những thao tác đó được thực hiện nhanh, chính xác và an toàn, các ứng dụng không tự quản lý dữ liệu trên các tệp tin rời rạc mà giao phó công việc này cho một \textbf{hệ quản trị cơ sở dữ liệu} (\emph{Database Management System} -- DBMS). Có thể hiểu DBMS là lớp phần mềm đứng giữa ứng dụng và dữ liệu: ứng dụng chỉ cần mô tả dữ liệu nào cần lưu hoặc cần lấy, còn việc tổ chức dữ liệu trên đĩa hay trên bộ nhớ, tìm ra cách truy xuất hiệu quả, giữ cho dữ liệu đúng khi nhiều người cùng thao tác và bảo vệ dữ liệu khi hệ thống gặp sự cố đều do DBMS đảm nhận. Nhờ vậy, người phát triển có thể tập trung vào nghiệp vụ của ứng dụng thay vì phải tự giải quyết những bài toán phức tạp về lưu trữ.

Tuy nhiên, không có một DBMS nào phù hợp cho mọi bài toán. Mỗi hệ quản trị cân bằng giữa các yêu cầu như tính toàn vẹn, khả năng truy vấn linh hoạt, tốc độ và độ bền của dữ liệu theo một cách riêng, và sự lựa chọn đó thể hiện rõ trong thiết kế bên trong của nó. Hiện nay có hai hướng tiếp cận tiêu biểu là các hệ quản trị cơ sở dữ liệu quan hệ (SQL) và các hệ NoSQL. Trong bài tập lớn này, nhóm chọn một đại diện cho mỗi hướng là \textbf{Microsoft SQL Server} và \textbf{Redis}. SQL Server là hệ quản trị cơ sở dữ liệu quan hệ được sử dụng rộng rãi trong doanh nghiệp: dữ liệu được tổ chức thành các bảng có lược đồ chặt chẽ, lưu trên đĩa và được truy vấn bằng ngôn ngữ T-SQL; bên cạnh đó, hệ quản trị này còn cung cấp hệ thống chỉ mục đa dạng, bộ tối ưu truy vấn, giao tác bảo đảm các tính chất ACID cùng cơ chế sao lưu và phục hồi đầy đủ. Ngược lại, Redis là một hệ lưu trữ khoá -- giá trị đặt toàn bộ dữ liệu trong bộ nhớ chính. Redis không có khái niệm lược đồ mà cung cấp các cấu trúc dữ liệu như chuỗi, bảng băm, danh sách, tập hợp và tập hợp có thứ tự, nhờ đó mỗi thao tác chỉ mất thời gian rất ngắn. Chính vì vậy, Redis thường được dùng làm bộ nhớ đệm, bộ đếm hoặc nơi lưu phiên đăng nhập. Có thể thấy hai hệ quản trị này đi theo hai triết lý gần như đối lập nhau: SQL Server ưu tiên sự toàn vẹn và độ bền của dữ liệu, trong khi Redis ưu tiên tốc độ. Cũng vì sự khác biệt đó mà trong thực tế, hai loại hệ quản trị này thường không thay thế nhau mà được kết hợp trong cùng một hệ thống, trong đó cơ sở dữ liệu quan hệ giữ vai trò nguồn dữ liệu gốc còn Redis đóng vai trò lớp phục vụ nhanh đứng phía trước.

Xuất phát từ những nhận định trên, bài tập lớn hướng tới hai mục tiêu gắn liền với nhau. Mục tiêu thứ nhất là tìm hiểu và giải thích các cơ chế bên trong của SQL Server và Redis, từ đó so sánh hai hệ quản trị một cách có cơ sở. Cụ thể, nhóm nghiên cứu năm chủ đề kỹ thuật, tương ứng với năm thành viên, gồm: quản lý và lưu trữ dữ liệu, cấu trúc chỉ mục, xử lý truy vấn, xử lý giao tác, sao lưu và phục hồi dữ liệu. Với mỗi chủ đề, nhóm trình bày cơ sở lý thuyết, phân tích cách mỗi hệ quản trị hiện thực vấn đề đó, rồi thiết kế các kịch bản thực nghiệm tương đương trên cả hai hệ để kiểm chứng và đối chiếu kết quả. Mục tiêu thứ hai là vận dụng những hiểu biết này vào việc xây dựng một ứng dụng thực tế. Ứng dụng được nhóm lựa chọn là \textbf{Novelle}, một trang web đọc và sáng tác truyện trực tuyến dành cho bốn nhóm người dùng: khách, người đọc, tác giả và quản trị viên. Đây là một bối cảnh phù hợp để kết hợp hai hệ quản trị, bởi vì dữ liệu của một nền tảng truyện vừa có nhiều quan hệ chặt chẽ cần được bảo đảm toàn vẹn, vừa có lượng truy cập đọc rất lớn cần được phục vụ nhanh. Trong Novelle, SQL Server lưu toàn bộ dữ liệu nghiệp vụ như tài khoản, truyện, chương và đánh giá, còn Redis đảm nhận bộ nhớ đệm, bộ đếm lượt xem, bản nháp, phiên đăng nhập và giới hạn tần suất. Cách phân chia vai trò này không được chọn theo cảm tính mà được giải thích dựa trên chính kết quả của các chủ đề nghiên cứu.

Để trình bày các nội dung trên một cách có hệ thống, phần còn lại của báo cáo được tổ chức như sau. Chương~\ref{sec:dbms-overview} giới thiệu tổng quan về kiến trúc và đặc điểm nổi bật của SQL Server và Redis, qua đó đưa ra một so sánh ban đầu giữa hai hệ quản trị. Tiếp theo, Chương~\ref{sec:app-db-design} mô tả ứng dụng Novelle cùng các yêu cầu chức năng, sau đó trình bày quá trình mô hình hoá dữ liệu, thiết kế lược đồ quan hệ trên SQL Server và cách tổ chức dữ liệu trên Redis. Chương~\ref{sec:dbms-topics} là phần trọng tâm của báo cáo, lần lượt trình bày năm chủ đề nghiên cứu, mỗi chủ đề gồm cơ sở lý thuyết, cơ chế trên từng hệ quản trị, thực nghiệm và đánh giá, và khép lại bằng phần kết luận chung cho cả hai hệ quản trị. Dựa trên những kết quả đó, Chương~\ref{sec:architecture-implementation} trình bày các công nghệ được sử dụng, kiến trúc hệ thống, cách tích hợp hai hệ quản trị và quá trình hiện thực các chức năng chính của ứng dụng. Cuối cùng, Chương~\ref{sec:conclusion} tổng kết những kết quả đã đạt được, đánh giá mức độ hoàn thành và chỉ ra những hạn chế còn tồn tại của đề tài. Danh mục tài liệu tham khảo và phụ lục chứa biên bản họp nhóm được đặt ở cuối báo cáo.

    \section{Tổng quan về hệ quản trị cơ sở dữ liệu}
\label{sec:dbms-overview}

Trước khi đi vào thiết kế ứng dụng và các chủ đề nghiên cứu chi tiết, chương này giới thiệu một cách tổng quan về hai hệ quản trị được sử dụng trong đề tài là Microsoft SQL Server và Redis. Với mỗi hệ quản trị, nhóm giới thiệu sơ lược về nguồn gốc, những đặc điểm nổi bật và các lĩnh vực ứng dụng phổ biến; phần cuối chương đặt hai hệ quản trị cạnh nhau để thấy rõ sự khác biệt trong cách tiếp cận của mỗi bên.

\subsection{Tổng quan về Microsoft SQL Server}
\label{subsec:mssql-overview}

Microsoft SQL Server là hệ quản trị cơ sở dữ liệu quan hệ do Microsoft phát triển, với phiên bản đầu tiên ra đời từ cuối những năm 1980. Trải qua hơn ba mươi năm phát triển, SQL Server đã trở thành một trong những hệ quản trị cơ sở dữ liệu phổ biến nhất trong môi trường doanh nghiệp, bên cạnh Oracle, MySQL và PostgreSQL. Ban đầu SQL Server chỉ chạy trên hệ điều hành Windows, nhưng hiện nay đã có thể chạy trên cả Linux, trong container và dưới dạng dịch vụ đám mây. Microsoft cung cấp nhiều phiên bản khác nhau, từ phiên bản Express miễn phí dành cho học tập và các ứng dụng nhỏ, đến các phiên bản Standard và Enterprise dành cho những hệ thống có quy mô lớn.

\begin{figure}[H]
    \centering
    \includegraphics[width=0.3\textwidth]{image/chapter2/sql-server.png}
    \caption{Microsoft SQL Server}
    \label{fig:sql-server-logo}
\end{figure}

\subsubsection{Đặc điểm nổi bật}

Là một hệ quản trị cơ sở dữ liệu quan hệ, SQL Server tổ chức dữ liệu thành các bảng có cấu trúc rõ ràng và liên kết với nhau thông qua các mối quan hệ. Người dùng làm việc với dữ liệu bằng ngôn ngữ SQL, một ngôn ngữ mang tính khai báo: chỉ cần mô tả dữ liệu mình muốn lấy, còn việc tìm ra cách lấy hiệu quả nhất do hệ quản trị tự đảm nhận. Nhờ vậy, những yêu cầu phức tạp như kết hợp thông tin từ nhiều bảng hay thống kê, tổng hợp số liệu đều có thể được thực hiện một cách linh hoạt.

Một đặc điểm quan trọng khác của SQL Server là khả năng bảo đảm \textbf{tính toàn vẹn} của dữ liệu. Hệ quản trị cho phép định nghĩa các quy tắc mà dữ liệu bắt buộc phải tuân theo và tự động từ chối những thay đổi vi phạm các quy tắc đó. Bên cạnh đó, SQL Server hỗ trợ đầy đủ các tính chất \textbf{ACID} của giao tác, nhờ đó một nhóm thao tác luôn được thực hiện trọn vẹn hoặc không được thực hiện gì cả, kể cả khi có lỗi hay có nhiều người dùng cùng thao tác. Dữ liệu được lưu bền vững trên đĩa và có thể được sao lưu, khôi phục khi gặp sự cố. Ngoài ra, SQL Server còn đi kèm một hệ sinh thái công cụ phong phú hỗ trợ quản trị, phân tích dữ liệu và lập báo cáo.

\subsubsection{Ứng dụng}

Với những đặc điểm trên, SQL Server thường được lựa chọn cho các hệ thống mà dữ liệu cần chính xác tuyệt đối và được lưu trữ lâu dài, chẳng hạn như hệ thống ngân hàng, kế toán, quản lý bán hàng, quản lý nhân sự hay các hệ thống hoạch định nguồn lực doanh nghiệp (ERP). Hệ quản trị này cũng thường được dùng làm nền tảng cho các kho dữ liệu và hệ thống báo cáo, phân tích kinh doanh. Nhìn chung, SQL Server phù hợp với những ứng dụng có dữ liệu nhiều quan hệ chặt chẽ và có nhu cầu truy vấn đa dạng.

\subsection{Tổng quan về Redis}
\label{subsec:redis-overview}

Redis (viết tắt của \emph{Remote Dictionary Server}) là một hệ lưu trữ dữ liệu thuộc nhóm NoSQL, được Salvatore Sanfilippo phát triển vào năm 2009. Ra đời từ nhu cầu cải thiện tốc độ cho một ứng dụng web, Redis nhanh chóng trở nên phổ biến nhờ tốc độ vượt trội và cách sử dụng đơn giản. Ngày nay, Redis được rất nhiều công ty công nghệ lớn sử dụng và có mặt trên hầu hết các nền tảng điện toán đám mây dưới dạng dịch vụ có sẵn.

\begin{figure}[H]
    \centering
    \includegraphics[width=0.3\textwidth]{image/chapter2/redis.png}
    \caption{Redis}
    \label{fig:redis-logo}
\end{figure}

\subsubsection{Đặc điểm nổi bật}

Khác với các hệ quản trị cơ sở dữ liệu truyền thống, Redis lưu toàn bộ dữ liệu trong \textbf{bộ nhớ chính} (RAM) thay vì trên đĩa. Đây chính là lý do giúp Redis có tốc độ đọc ghi rất cao, mỗi thao tác thường chỉ mất một khoảng thời gian rất nhỏ. Redis tổ chức dữ liệu theo mô hình \textbf{khoá -- giá trị}: mỗi dữ liệu được gắn với một khoá duy nhất và được truy cập trực tiếp thông qua khoá đó. Điểm đặc biệt là giá trị không chỉ là một chuỗi đơn giản mà có thể là nhiều dạng cấu trúc dữ liệu khác nhau như danh sách, tập hợp hay tập hợp có thứ tự, giúp người phát triển dễ dàng biểu diễn những dữ liệu như bảng xếp hạng, bộ đếm hay hàng đợi.

Bên cạnh tốc độ, Redis còn có cách sử dụng đơn giản và linh hoạt, vì không đòi hỏi phải định nghĩa trước cấu trúc của dữ liệu. Mỗi dữ liệu có thể được đặt thời gian sống và tự động bị xoá khi hết hạn, điều rất thuận tiện cho các dữ liệu tạm thời. Mặc dù làm việc chủ yếu trên bộ nhớ, Redis vẫn có thể ghi dữ liệu xuống đĩa để khôi phục khi khởi động lại. Tuy nhiên, đổi lại cho tốc độ và sự đơn giản, Redis không hỗ trợ truy vấn phức tạp như SQL, không có các ràng buộc toàn vẹn, và lượng dữ liệu có thể lưu bị giới hạn bởi dung lượng bộ nhớ.

\subsubsection{Ứng dụng}

Nhờ tốc độ cao, Redis thường được dùng làm \textbf{bộ nhớ đệm} (cache) đứng trước một cơ sở dữ liệu chính, giúp giảm tải và rút ngắn thời gian phản hồi cho những dữ liệu được truy cập thường xuyên. Ngoài ra, Redis còn được ứng dụng rộng rãi để lưu phiên đăng nhập của người dùng, xây dựng bảng xếp hạng trong trò chơi, đếm lượt xem và lượt thích, giới hạn số lần truy cập của người dùng, làm hàng đợi công việc hoặc truyền tin nhắn theo thời gian thực. Nói cách khác, Redis phù hợp với những dữ liệu cần được truy cập tức thời, có cấu trúc đơn giản và có thể chấp nhận mất một phần nhỏ khi xảy ra sự cố.

\subsection{So sánh tổng quan MS SQL Server và Redis}
\label{subsec:mssql-vs-redis}

Từ những giới thiệu trên, có thể thấy SQL Server và Redis được thiết kế cho những mục đích khá khác nhau. Bảng~\ref{tab:overview-comparison} tóm tắt sự khác biệt giữa hai hệ quản trị ở một số khía cạnh chính.

\begin{table}[H]
\centering
\small
\renewcommand{\arraystretch}{1.3}
\begin{tabularx}{\textwidth}{|>{\raggedright\arraybackslash}p{3cm}|>{\raggedright\arraybackslash}X|>{\raggedright\arraybackslash}X|}
\hline
\textbf{Tiêu chí} & \textbf{Microsoft SQL Server} & \textbf{Redis} \\ \hline
Loại hệ quản trị & Hệ quản trị cơ sở dữ liệu quan hệ (SQL). & Hệ lưu trữ khoá -- giá trị trong bộ nhớ (NoSQL). \\ \hline
Cách tổ chức dữ liệu & Các bảng có cấu trúc định trước và quan hệ với nhau. & Các cặp khoá -- giá trị, không cần định nghĩa cấu trúc trước. \\ \hline
Nơi lưu dữ liệu & Chủ yếu trên đĩa. & Chủ yếu trên bộ nhớ chính. \\ \hline
Thế mạnh & Dữ liệu toàn vẹn, an toàn; truy vấn linh hoạt và phức tạp. & Tốc độ đọc ghi rất cao; sử dụng đơn giản. \\ \hline
Hạn chế & Tốc độ chậm hơn với lượng truy cập rất lớn. & Không hỗ trợ truy vấn phức tạp; dung lượng bị giới hạn bởi bộ nhớ; có thể mất một ít dữ liệu khi gặp sự cố. \\ \hline
Ứng dụng tiêu biểu & Ngân hàng, kế toán, quản lý bán hàng, hệ thống doanh nghiệp. & Bộ nhớ đệm, phiên đăng nhập, bộ đếm, bảng xếp hạng, hàng đợi. \\ \hline
\end{tabularx}
\caption{So sánh tổng quan SQL Server và Redis}
\label{tab:overview-comparison}
\end{table}

Qua bảng so sánh, có thể thấy hai hệ quản trị không cạnh tranh trực tiếp với nhau mà mỗi bên mạnh ở một khía cạnh riêng: SQL Server đặt sự an toàn và toàn vẹn của dữ liệu lên hàng đầu, còn Redis ưu tiên tốc độ. Chính vì vậy, trong nhiều hệ thống thực tế, hai loại hệ quản trị này thường được sử dụng kết hợp với nhau. Những khác biệt này sẽ được phân tích sâu hơn về mặt cơ chế và kiểm chứng bằng thực nghiệm trong Chương~\ref{sec:dbms-topics}.

    \section{Tổng quan ứng dụng và thiết kế cơ sở dữ liệu}
\label{sec:app-db-design}

% ===================================================================
\subsection{Tổng quan ứng dụng}
\label{subsec:app-overview}

\subsubsection{Giới thiệu ứng dụng}

\textbf{Novelle} là một ứng dụng web cho phép người dùng đọc và sáng tác truyện dài kỳ trực tuyến. Nội dung truyện do chính người dùng viết và đăng tải theo từng chương. Ứng dụng được xây dựng theo mô hình client--server: phía client là một ứng dụng web chạy trên trình duyệt, phía server cung cấp REST API để xử lý nghiệp vụ và truy xuất dữ liệu.

\subsubsection{Mục tiêu và phạm vi ứng dụng}

Ứng dụng hướng tới việc cung cấp đầy đủ các luồng chính của một nền tảng truyện trực tuyến, bao gồm: tìm và đọc truyện, quản lý thư viện cá nhân, sáng tác và phát hành chương, và quản trị nội dung. Phạm vi của ứng dụng giới hạn ở các chức năng cốt lõi này, trong đó ưu tiên các thao tác thêm, sửa, xoá và truy vấn dữ liệu được hiện thực đầy đủ.

\subsubsection{Các chức năng chính của ứng dụng}

\begin{itemize}
    \item \textbf{Trang chủ và tìm kiếm:} hiển thị truyện mới cập nhật, xem nhiều, đánh giá cao; tìm truyện theo từ khoá và lọc theo thể loại, thẻ, trạng thái, điểm đánh giá.
    \item \textbf{Chi tiết truyện và đọc chương:} xem thông tin truyện, mục lục, trang tác giả; đọc chương và tự động lưu vị trí đọc.
    \item \textbf{Thư viện cá nhân:} lịch sử đọc, truyện và tác giả đang theo dõi, tủ truyện tự tạo.
    \item \textbf{Đánh giá:} chấm điểm và viết nhận xét cho truyện.
    \item \textbf{Đăng ký làm tác giả:} gửi yêu cầu kèm bút danh và theo dõi trạng thái xét duyệt.
    \item \textbf{Studio sáng tác:} tạo và quản lý truyện, viết chương với bản nháp tự động lưu, chia tập, hẹn giờ đăng, xem lại các phiên bản chỉnh sửa, và xem thống kê lượt xem, lượt theo dõi, điểm đánh giá.
    \item \textbf{Quản trị:} quản lý và khoá tài khoản, duyệt yêu cầu tác giả, ẩn hoặc gỡ truyện và nhận xét vi phạm, quản lý thể loại và thẻ, xem thống kê toàn hệ thống.
\end{itemize}

% ===================================================================
\subsection{Yêu cầu chức năng của hệ thống}
\label{subsec:functional-requirements}

\subsubsection{Danh sách các yêu cầu chức năng}
% TODO

\subsubsection{Mô tả chi tiết các chức năng}
% TODO

% ===================================================================
\subsection{Phân tích và mô hình hóa dữ liệu}
\label{subsec:data-modeling}

\subsubsection{Xác định các thực thể dữ liệu}
% TODO

\subsubsection{Xác định thuộc tính của các thực thể}
% TODO

\subsubsection{Xác định mối quan hệ giữa các thực thể}
% TODO

\subsubsection{Xây dựng sơ đồ thực thể - liên kết (ERD)}
\label{subsubsec:erd}
% TODO: chèn sơ đồ ERD, ví dụ:
% \begin{figure}[H]
%     \centering
%     \includegraphics[width=\textwidth]{image/chapter3/erd.png}
%     \caption{Sơ đồ thực thể - liên kết của hệ thống}
%     \label{fig:erd}
% \end{figure}

% ===================================================================
\subsection{Thiết kế cơ sở dữ liệu trên MS SQL Server (SQL)}
\label{subsec:mssql-design}

\subsubsection{Xây dựng lược đồ quan hệ (Relational Schema)}
% TODO

\subsubsection{Thiết kế các bảng và thuộc tính}
% TODO

\subsubsection{Xác định khóa chính và khóa ngoại}
% TODO

\subsubsection{Xác định các ràng buộc dữ liệu}
% TODO

% ===================================================================
\subsection{Thiết kế dữ liệu trên Redis (NoSQL)}
\label{subsec:redis-design}

\subsubsection{Chiến lược tổ chức dữ liệu}
% TODO

\subsubsection{Thiết kế Key Naming Convention}
% TODO

\subsubsection{Lựa chọn cấu trúc dữ liệu Redis}
% TODO

\subsubsection{Thiết kế cách lưu trữ và truy xuất dữ liệu}
% TODO

    \section{So sánh và nghiên cứu}
\label{sec:dbms-topics}

Đây là chương trọng tâm của báo cáo. Trong chương này, nhóm lần lượt trình bày năm chủ đề nghiên cứu: quản lý và lưu trữ dữ liệu, cấu trúc chỉ mục, xử lý truy vấn, xử lý giao tác, sao lưu và phục hồi dữ liệu. Mỗi chủ đề được trình bày theo cùng một trình tự: trước hết là cơ sở lý thuyết, tiếp đến là cơ chế mà SQL Server và Redis sử dụng để giải quyết vấn đề, sau đó là các kịch bản thực nghiệm được thiết kế trên chính dữ liệu và nghiệp vụ của Novelle, và cuối cùng là phần so sánh, đánh giá. Các thực nghiệm đều được chạy bằng công cụ lab do nhóm tự xây dựng, trong đó SQL Server chạy trên AWS RDS còn Redis chạy trong Docker trên máy cá nhân. Vì hai hệ quản trị không chạy trên cùng một hạ tầng, nhóm chủ yếu so sánh thời gian xử lý phía server, tính đúng của kết quả và hành vi của mỗi hệ, thay vì so sánh trực tiếp độ trễ mà người dùng cảm nhận. Chương khép lại bằng phần kết luận chung, tổng hợp những gì rút ra được từ cả năm chủ đề.

\subsection{Quản lý và lưu trữ dữ liệu (Data Storage and Management)}
\label{subsec:data-storage}

\subsubsection{Cơ sở lý thuyết}
% TODO

\subsubsection{Cơ chế lưu trữ trên MS SQL Server}
% TODO

\subsubsection{Cơ chế lưu trữ trên Redis}
% TODO

\subsubsection{Thiết kế kịch bản và thực hiện thực nghiệm}
% TODO

\subsubsection{So sánh và đánh giá}
% TODO

\subsection{Cấu trúc chỉ mục (Indexing)}
\label{subsec:indexing}
\subsubsection{Cơ sở lý thuyết}
\label{subsubsec:indexing-theory}

Chỉ mục (\emph{index}) là một cấu trúc dữ liệu phụ được xây dựng trên một hoặc nhiều thuộc tính, gọi là \emph{khoá tìm kiếm}, nhằm giúp hệ quản trị tìm ra các bản ghi thoả điều kiện mà không phải quét qua toàn bộ tệp dữ liệu. Có thể hình dung chỉ mục giống như mục lục ở cuối một cuốn sách: thay vì lật từng trang để tìm một thuật ngữ, người đọc chỉ cần tra mục lục để biết ngay thuật ngữ đó nằm ở trang nào.

Dựa vào mối liên hệ với thứ tự vật lý của dữ liệu, chỉ mục được chia thành nhiều loại. \textbf{Chỉ mục chính} và \textbf{chỉ mục gom cụm} gắn liền với thứ tự sắp xếp vật lý của tệp, vì vậy mỗi tệp chỉ có tối đa một chỉ mục loại này. Ngược lại, \textbf{chỉ mục phụ} có thể được xây dựng trên bất kỳ thuộc tính nào và một tệp có thể có nhiều chỉ mục phụ. Bên cạnh đó, chỉ mục còn được phân biệt thành \textbf{chỉ mục dày}, có một mục cho mỗi bản ghi, và \textbf{chỉ mục thưa}, chỉ có một mục cho mỗi khối dữ liệu. Khi khoá tìm kiếm gồm nhiều thuộc tính, ta có \textbf{chỉ mục kết hợp}; với loại chỉ mục này, thứ tự của các thuộc tính sẽ quyết định những điều kiện truy vấn nào có thể tận dụng được chỉ mục.

Về mặt cấu trúc, có ba dạng chỉ mục phổ biến. \textbf{Cây B\textsuperscript{+}} là cây cân bằng nhiều cấp, cho phép tìm kiếm với chi phí $O(\log N)$; do các nút lá được liên kết với nhau theo thứ tự, cây B\textsuperscript{+} hỗ trợ rất tốt các truy vấn theo khoảng và các truy vấn cần sắp xếp. \textbf{Chỉ mục băm} cho phép tìm kiếm theo điều kiện bằng với chi phí chỉ khoảng $O(1)$, nhưng lại không hỗ trợ truy vấn theo khoảng. Cuối cùng, \textbf{chỉ mục đảo} ánh xạ mỗi từ tới danh sách các tài liệu chứa từ đó, và thường được dùng cho chức năng tìm kiếm toàn văn.

Tuy nhiên, chỉ mục không phải là miễn phí. Chỉ mục giúp việc đọc nhanh hơn, nhưng đồng thời cũng tốn thêm dung lượng lưu trữ và làm chậm các thao tác ghi, bởi vì mỗi lệnh \texttt{INSERT}, \texttt{UPDATE} hay \texttt{DELETE} đều phải cập nhật cả các chỉ mục liên quan. Vì vậy, việc thiết kế chỉ mục luôn là bài toán cân bằng giữa tốc độ đọc và chi phí ghi.

% -------------------------------------------------------------------
\subsubsection{Cơ chế chỉ mục trên MS SQL Server}
\label{subsubsec:indexing-mssql}

Với nền tảng lý thuyết trên, phần này tìm hiểu cách SQL Server tổ chức chỉ mục trên đĩa và các loại chỉ mục mà hệ quản trị này cung cấp.

\paragraph{Tổ chức lưu trữ nền tảng.} SQL Server lưu dữ liệu trong các trang (\emph{page}) kích thước 8~KB, nhóm thành các extent gồm 8 trang liên tiếp. Chỉ mục \emph{rowstore} của SQL Server được tổ chức dưới dạng cây B\textsuperscript{+}: nút gốc và các nút trung gian chứa khoá dẫn đường, nút lá chứa dữ liệu chỉ mục; các trang ở cùng một cấp được liên kết đôi với nhau để duyệt xuôi hoặc ngược theo thứ tự khoá. Một bảng có thể được lưu ở một trong hai dạng:

\begin{itemize}
    \item \textbf{Heap:} bảng không có chỉ mục gom cụm; các dòng không theo thứ tự nào và được định vị bằng \emph{RID} (mã tệp, mã trang, số khe trong trang).
    \item \textbf{Bảng có chỉ mục gom cụm:} bản thân dữ liệu được sắp xếp và lưu trong các nút lá của cây B\textsuperscript{+}.
\end{itemize}

\paragraph{Clustered index.} Chỉ mục gom cụm (\emph{clustered index}) quyết định thứ tự lưu trữ vật lý của các dòng; nút lá của nó \emph{chính là} các trang dữ liệu, nên mỗi bảng chỉ có tối đa một clustered index. Mặc định, khi khai báo \texttt{PRIMARY KEY}, SQL Server tạo một clustered index trên khoá chính. Khoá gom cụm nên ngắn, duy nhất, ít thay đổi và tăng dần, vì nó được chép vào mọi chỉ mục phụ và vì khoá tăng dần giúp các dòng mới luôn được ghi vào cuối cây, hạn chế tách trang.

\paragraph{Nonclustered index.} Chỉ mục không gom cụm (\emph{nonclustered index}) là một cây B\textsuperscript{+} riêng, tách khỏi dữ liệu, tương ứng với chỉ mục phụ trong lý thuyết. Mỗi dòng ở nút lá gồm giá trị khoá chỉ mục và một \emph{bộ định vị dòng} (\emph{row locator}): là RID nếu bảng là heap, hoặc là giá trị khoá gom cụm nếu bảng có clustered index. Một bảng có thể có tới 999 nonclustered index. Khi truy vấn cần thêm các cột không có trong chỉ mục, SQL Server phải dùng bộ định vị để tra lại dòng đầy đủ, thao tác này gọi là \emph{Key Lookup} (hoặc \emph{RID Lookup} với heap). Vì mỗi lần tra là một lần truy cập ngẫu nhiên, khi số dòng khớp vượt một ngưỡng (\emph{tipping point}), bộ tối ưu truy vấn sẽ bỏ qua chỉ mục và chọn quét toàn bảng.

\paragraph{Các biến thể của nonclustered index.} SQL Server hỗ trợ nhiều tuỳ chọn để làm chỉ mục phù hợp hơn với truy vấn:

\begin{itemize}
    \item \textbf{Chỉ mục kết hợp} (\emph{composite}): khoá gồm tối đa 32 cột. Chỉ mục được tìm kiếm hiệu quả (\emph{Index Seek}) khi điều kiện truy vấn dùng các cột theo thứ tự từ trái sang; thông thường các cột so sánh bằng nên đứng trước cột so sánh theo khoảng.
    \item \textbf{Chỉ mục bao phủ} (\emph{covering index}) với mệnh đề \texttt{INCLUDE}: các cột được \texttt{INCLUDE} chỉ nằm ở nút lá, không tham gia sắp xếp, giúp chỉ mục chứa đủ mọi cột truy vấn cần và loại bỏ hoàn toàn bước Key Lookup.
    \item \textbf{Chỉ mục duy nhất} (\emph{unique index}): bảo đảm không có hai dòng trùng giá trị khoá. Ràng buộc \texttt{UNIQUE} được SQL Server hiện thực bằng chính loại chỉ mục này.
    \item \textbf{Chỉ mục có lọc} (\emph{filtered index}): chỉ đánh chỉ mục những dòng thoả mệnh đề \texttt{WHERE}, nên nhỏ hơn, rẻ hơn khi ghi và có thể kết hợp với \texttt{UNIQUE} để biểu diễn ràng buộc chỉ áp dụng cho một phần dữ liệu.
    \item \textbf{Thứ tự khoá} \texttt{ASC}/\texttt{DESC} cho từng cột, giúp tránh bước sắp xếp khi truy vấn có \texttt{ORDER BY} khớp với thứ tự của chỉ mục.
\end{itemize}

\paragraph{Chỉ mục toàn văn (Full-Text Index).} Để tìm kiếm theo từ trong văn bản dài, SQL Server cung cấp \emph{full-text index}, một chỉ mục đảo được lưu trong \emph{full-text catalog}. Khi đánh chỉ mục, văn bản được tách từ bởi \emph{word breaker} của ngôn ngữ tương ứng, và có thể bỏ dấu khi catalog được tạo với \texttt{ACCENT\_SENSITIVITY = OFF}. Truy vấn dùng các vị từ \texttt{CONTAINS}, \texttt{FREETEXT} hoặc các hàm \texttt{CONTAINSTABLE}, \texttt{FREETEXTTABLE} để hỗ trợ tìm theo tiền tố, cụm từ và xếp hạng mức độ liên quan. Mỗi bảng chỉ có một full-text index.

\paragraph{Các loại chỉ mục khác.} SQL Server còn có \emph{columnstore index} (lưu dữ liệu theo cột, nén cao, phù hợp truy vấn phân tích trên bảng lớn), \emph{spatial index}, \emph{XML index} và chỉ mục băm cho bảng tối ưu bộ nhớ (\emph{memory-optimized}).

\paragraph{Thống kê và bộ tối ưu truy vấn.} SQL Server không bắt buộc người dùng chỉ định chỉ mục. Với mỗi chỉ mục, hệ thống duy trì \emph{thống kê} (\emph{statistics}) gồm histogram phân bố giá trị của cột khoá đầu tiên và mật độ của các tiền tố khoá. Bộ tối ưu dựa trên chi phí (\emph{cost-based optimizer}) dùng thống kê này để ước lượng số dòng khớp và tự quyết định giữa Index Seek, Index Scan, Key Lookup hay Table Scan. Thống kê được tự động cập nhật khi lượng dữ liệu thay đổi vượt một ngưỡng.

\paragraph{Bảo trì chỉ mục.} Khi một dòng được chèn vào trang đã đầy, SQL Server tách trang (\emph{page split}), chuyển một nửa số dòng sang trang mới. Theo thời gian, việc tách trang gây \emph{phân mảnh} (\emph{fragmentation}): thứ tự logical của các trang không còn trùng với thứ tự vật lý, và các trang chứa nhiều khoảng trống. Tham số \texttt{FILLFACTOR} cho phép chừa sẵn khoảng trống khi tạo chỉ mục, còn các lệnh \texttt{ALTER INDEX \ldots{} REORGANIZE} và \texttt{ALTER INDEX \ldots{} REBUILD} dùng để sắp xếp lại hoặc xây lại chỉ mục khi phân mảnh tăng cao.

% -------------------------------------------------------------------
\subsubsection{Cơ chế chỉ mục trên Redis}
\label{subsubsec:indexing-redis}

Trái ngược với sự đa dạng của SQL Server, Redis lõi hầu như không có khái niệm chỉ mục theo nghĩa truyền thống. Phần này giải thích vì sao lại như vậy và các cách mà người phát triển có thể dùng để xây dựng chỉ mục trên Redis.

\paragraph{Chỉ mục chính: không gian khoá là một bảng băm.} Redis là kho key--value lưu toàn bộ dữ liệu trong bộ nhớ. Mỗi cơ sở dữ liệu logic của Redis là một từ điển (\emph{dict}) ánh xạ tên key tới giá trị, được hiện thực bằng bảng băm với kỹ thuật xử lý xung đột bằng danh sách liên kết. Vì vậy, truy cập theo tên key luôn có chi phí trung bình $O(1)$ và có thể xem tên key là ``chỉ mục chính'' duy nhất mà Redis cung cấp sẵn. Khi bảng băm cần mở rộng hoặc thu nhỏ, Redis không băm lại toàn bộ cùng lúc mà thực hiện \emph{băm lại tăng dần} (\emph{incremental rehashing}): duy trì song song hai bảng và di chuyển dần từng bucket sang bảng mới qua mỗi thao tác, để không có thao tác nào bị chặn lâu.

\paragraph{Không có chỉ mục phụ tự động.} Redis lõi không có khái niệm lược đồ, cột hay chỉ mục phụ. Muốn tìm các giá trị theo một trường không phải tên key (ví dụ ``các sản phẩm thuộc một danh mục có giá dưới 100''), có hai cách: duyệt toàn bộ không gian khoá bằng \texttt{SCAN} với chi phí $O(N)$, hoặc để ứng dụng \textbf{tự xây dựng và tự duy trì} cấu trúc chỉ mục bằng các kiểu dữ liệu sẵn có. Mọi lệnh ghi dữ liệu gốc khi đó phải đi kèm lệnh cập nhật chỉ mục, thường được gộp trong \texttt{MULTI}/\texttt{EXEC} hoặc một script Lua để giữ tính nhất quán.

\paragraph{Chỉ mục phụ tự xây dựng bằng kiểu dữ liệu Redis.} Các mẫu chỉ mục phổ biến gồm:

\begin{itemize}
    \item \textbf{Set làm chỉ mục theo giá trị:} mỗi giá trị của một trường có một Set chứa các định danh tương ứng, ví dụ \texttt{category:book} chứa id các sản phẩm thuộc danh mục sách. Điều kiện ``và'' giữa nhiều trường được thực hiện bằng phép giao \texttt{SINTER}, điều kiện ``hoặc'' bằng phép hợp \texttt{SUNION}.
    \item \textbf{Sorted Set làm chỉ mục theo khoảng và sắp xếp:} điểm (\emph{score}) của mỗi phần tử là giá trị số cần sắp xếp, ví dụ giá bán hoặc thời điểm tạo dạng epoch. Sorted Set được hiện thực bằng một \emph{skip list} kết hợp với bảng băm, nên tìm theo khoảng điểm (\texttt{ZRANGE \ldots{} BYSCORE}) có chi phí $O(\log N + M)$ và lấy theo thứ hạng (\texttt{ZRANGE} với vị trí) cũng có chi phí $O(\log N + M)$, với $M$ là số phần tử trả về. Với tập nhỏ, Redis dùng biểu diễn \emph{listpack} gọn hơn để tiết kiệm bộ nhớ.
    \item \textbf{Sorted Set làm chỉ mục kết hợp:} nhúng trường so sánh bằng vào tên key và dùng trường khoảng làm score, ví dụ \texttt{by\_price:book}, tương tự một chỉ mục kết hợp (\texttt{Category}, \texttt{Price}) của SQL Server. Khi mọi phần tử có cùng score, Sorted Set được sắp theo thứ tự từ điển, cho phép tìm theo tiền tố chuỗi bằng \texttt{ZRANGE \ldots{} BYLEX}.
\end{itemize}

Ví dụ một chỉ mục phụ tự xây dựng:

\begin{minted}[fontsize=\small,frame=single,breaklines]{text}
HSET product:42 Name "Notebook" Category book Price 45
SADD category:book 42                         -- chỉ mục theo giá trị
ZADD by_price 45 42                           -- chỉ mục theo khoảng
ZRANGE by_price 0 100 BYSCORE                 -- sản phẩm có giá <= 100
\end{minted}

\paragraph{Module RediSearch: chỉ mục phụ do Redis quản lý.} Redis Stack tích hợp module \textbf{RediSearch}, bổ sung khả năng đánh chỉ mục phụ tự động cho dữ liệu dạng Hash hoặc JSON. Lệnh \texttt{FT.CREATE} khai báo một chỉ mục gồm tiền tố key cần theo dõi, một lược đồ các trường cùng loại chỉ mục, và tuỳ chọn một biểu thức lọc. Sau đó, mọi lệnh ghi lên các Hash có tiền tố tương ứng đều được module tự động cập nhật vào chỉ mục, ứng dụng không phải tự duy trì. Các loại trường chính gồm:

\begin{itemize}
    \item \texttt{TEXT}: chỉ mục đảo dùng cho tìm kiếm toàn văn, có tách từ, xếp hạng mức độ liên quan và tuỳ chọn \texttt{NOSTEM} để tắt việc rút gọn từ về gốc.
    \item \texttt{TAG}: chỉ mục theo giá trị chính xác, phù hợp với trạng thái, thể loại hoặc các mã liệt kê.
    \item \texttt{NUMERIC}: chỉ mục theo khoảng số, được tổ chức dưới dạng cây phân khoảng.
    \item \texttt{GEO} và \texttt{VECTOR}: chỉ mục toạ độ địa lý và chỉ mục vector cho tìm kiếm tương đồng.
\end{itemize}

Tuỳ chọn \texttt{SORTABLE} giữ thêm một bản sao giá trị để sắp xếp kết quả nhanh với \texttt{SORTBY}, còn \texttt{FILTER} cho phép chỉ đánh chỉ mục các Hash thoả một điều kiện, tương tự filtered index của SQL Server. Truy vấn được thực hiện bằng \texttt{FT.SEARCH} hoặc \texttt{FT.AGGREGATE}:

\begin{minted}[fontsize=\small,frame=single,breaklines]{text}
FT.CREATE idx:product ON HASH PREFIX 1 product:
    SCHEMA Name TEXT NOSTEM
           Category TAG
           Price NUMERIC SORTABLE
FT.SEARCH idx:product "@Category:{book} @Price:[0 100]"
    SORTBY Price ASC LIMIT 0 20
\end{minted}

\paragraph{Đặc điểm chung của chỉ mục trên Redis.} Dù xây dựng thủ công hay qua RediSearch, chỉ mục trên Redis đều nằm hoàn toàn trong bộ nhớ, nên mỗi chỉ mục làm tăng trực tiếp lượng RAM sử dụng. Redis cũng không có bộ tối ưu truy vấn dựa trên chi phí như SQL Server: với chỉ mục tự xây dựng, ứng dụng phải tự chọn cấu trúc và thứ tự thao tác; với RediSearch, kế hoạch thực thi được suy ra trực tiếp từ biểu thức truy vấn. Ngoài ra, \texttt{FT.SEARCH} mặc định chỉ trả về tối đa 10\,000 kết quả (\texttt{MAXSEARCHRESULTS}), nên việc phân trang sâu trên tập kết quả lớn bị giới hạn.

% -------------------------------------------------------------------
\subsubsection{Thiết kế kịch bản và thực hiện thực nghiệm}
\label{subsubsec:indexing-experiment}

% -------------------------------------------------------------------
\paragraph{Kịch bản 1 -- Tra cứu theo một cột.}

\textbf{Mục đích.} Kịch bản này nhằm kiểm tra lợi ích cơ bản nhất của chỉ mục là tìm nhanh các dòng theo giá trị của một cột không phải khoá chính. Hai truy vấn được dùng: tìm đúng một truyện theo \texttt{Slug} (giá trị duy nhất) và tìm các truyện của một tác giả theo \texttt{AuthorId} (vài dòng).

\textbf{SQL Server:} chạy truy vấn khi chưa có chỉ mục (phải quét toàn bảng) và sau khi tạo một nonclustered index trên cột cần tìm.

\begin{minted}[fontsize=\small,frame=single,breaklines]{sql}
SELECT StoryId, Title, Slug, CompletionStatus, ChapterCount, RatingAvg
FROM Story WHERE Slug = @slug;

SELECT StoryId, Title, Slug, CompletionStatus, ChapterCount, RatingAvg
FROM Story WHERE AuthorId = @authorId

CREATE NONCLUSTERED INDEX IX_Slug   ON Story (Slug);
CREATE NONCLUSTERED INDEX IX_Author ON Story (AuthorId);
\end{minted}

\textbf{Redis:} so sánh ba cách: duyệt toàn bộ key bằng \texttt{SCAN} rồi lọc từng Hash; chỉ mục tự xây (String ánh xạ slug sang id, Set chứa id truyện của mỗi tác giả); và chỉ mục RediSearch với trường \texttt{TAG}.

\begin{minted}[fontsize=\small,frame=single,breaklines]{text}
SCAN 0 MATCH story:* COUNT 1000          -- rồi HGET từng Hash để so sánh Slug
GET slug:<slug>        -> HGETALL story:<id>       -- index tự xây
SMEMBERS author:<id>   -> HGETALL story:<id> ...
FT.SEARCH idx:story "@Slug:{<slug>}"                -- RediSearch TAG
\end{minted}

\begin{table}[H]
\centering\small
\begin{tabular}{|l|l|r|r|}
\hline
\textbf{DBMS} & \textbf{Biến thể} & \textbf{Theo Slug (ms)} & \textbf{Theo AuthorId (ms)} \\ \hline
SQL Server & Không chỉ mục (đọc 74 trang) & 1,665 & 0,716 \\ \hline
SQL Server & Nonclustered index (đọc 4 / 10 trang) & 0,054 & 0,053 \\ \hline
Redis & \texttt{SCAN} toàn bộ key & 2,383 & 1,923 \\ \hline
Redis & Chỉ mục tự xây & 0,004 & 0,032 \\ \hline
Redis & RediSearch \texttt{TAG} & 0,029 & 0,039 \\ \hline
\end{tabular}
\caption{Kết quả kịch bản 1 (1\,000 truyện)}
\label{tab:idx-i1}
\end{table}

\textbf{Nhận xét.} Qua kết quả trên, có thể thấy chỉ mục giúp SQL Server đọc từ 74 trang xuống còn 4--10 trang và nhanh hơn 13--30 lần. Redis không có chỉ mục phụ sẵn nên khi dùng \texttt{SCAN} còn chậm hơn SQL Server quét bảng; khi có chỉ mục tự xây, việc tra cứu chỉ còn một lệnh \texttt{GET} và nhanh nhất trong các biến thể.

% -------------------------------------------------------------------
\paragraph{Kịch bản 2 -- Chỉ mục nhiều cột.}

\textbf{Mục đích.} Kịch bản này nhằm kiểm tra chỉ mục khi truy vấn lọc theo nhiều điều kiện, như bộ lọc ``truyện đã hoàn thành có điểm đánh giá từ X trở lên'' (khoảng 5\% số truyện khớp). Kịch bản so sánh chỉ mục kết hợp (composite) với chỉ mục bao phủ (covering), loại chứa sẵn cả các cột cần trả về.

\textbf{SQL Server:}

\begin{minted}[fontsize=\small,frame=single,breaklines]{sql}
SELECT StoryId, Title, RatingAvg, ChapterCount
FROM Story
WHERE CompletionStatus = @status AND RatingAvg >= @minRating;

-- composite: chỉ chứa 2 cột lọc
CREATE NONCLUSTERED INDEX IX_Status_Rating ON Story (CompletionStatus, RatingAvg);
-- covering: thêm các cột cần trả về bằng INCLUDE
CREATE NONCLUSTERED INDEX IX_Status_Rating_Cover ON Story (CompletionStatus, RatingAvg)
    INCLUDE (Title, ChapterCount);
\end{minted}

\textbf{Redis:} ngoài \texttt{SCAN}, dùng một Sorted Set cho mỗi trạng thái với điểm là \texttt{RatingAvg} (tên key mang điều kiện ``bằng'', điểm mang điều kiện ``khoảng'', tương đương chỉ mục kết hợp), hoặc RediSearch với trường \texttt{TAG} và \texttt{NUMERIC}.

\begin{minted}[fontsize=\small,frame=single,breaklines]{text}
ZRANGE by_rating:<status> <minRating> +inf BYSCORE   -- rồi HMGET từng truyện
FT.SEARCH idx:story "@CompletionStatus:{<status>} @RatingAvg:[<minRating> +inf]"
\end{minted}

\begin{table}[H]
\centering\small
\begin{tabular}{|l|l|r|r|r|}
\hline
\textbf{DBMS} & \textbf{Biến thể} & \textbf{100 truyện (ms)} & \textbf{1\,000 truyện (ms)} & \textbf{Trang đọc (1\,000)} \\ \hline
SQL Server & Không chỉ mục & 0,116 & 1,137 & 74 \\ \hline
SQL Server & Composite & 0,114 & 1,104 & 74 \\ \hline
SQL Server & Covering & 0,125 & 0,515 & 5 \\ \hline
Redis & \texttt{SCAN} toàn bộ key & 0,465 & 2,224 & -- \\ \hline
Redis & Sorted Set theo trạng thái & 0,036 & 0,150 & -- \\ \hline
Redis & RediSearch & 0,048 & 0,158 & -- \\ \hline
\end{tabular}
\caption{Kết quả kịch bản 2}
\label{tab:idx-i2}
\end{table}

\textbf{Nhận xét.} Kết quả cho thấy dù có chỉ mục phức hợp, bộ tối ưu của SQL Server vẫn chọn quét bảng (74 trang), vì với khoảng 50 dòng khớp, việc tra ngược từng dòng để lấy \texttt{Title}, \texttt{ChapterCount} (Key Lookup) tốn hơn quét toàn bộ. Chỉ mục bao phủ chứa sẵn các cột này nên được dùng, chỉ đọc 5 trang và nhanh gấp đôi. Redis không có bộ tối ưu: cấu trúc nào được dùng là do ứng dụng chọn, và cả Sorted Set lẫn RediSearch đều nhanh hơn \texttt{SCAN} khoảng 14 lần.

% -------------------------------------------------------------------
\paragraph{Kịch bản 3 -- Sắp xếp và phân trang.}

\textbf{Mục đích.} Kịch bản này nhằm kiểm tra chỉ mục khi truy vấn cần kết quả đã sắp xếp và chia trang, như danh sách ``truyện mới cập nhật'' trên trang chủ, mỗi trang 20 truyện. Mỗi biến thể được đo ở trang đầu và ở trang nằm giữa danh sách, để thấy ảnh hưởng của độ sâu trang. Hai cách phân trang được so sánh: \texttt{OFFSET} (bỏ qua $n$ dòng đầu) và \emph{keyset} (lấy tiếp các dòng sau dòng cuối của trang trước).

\textbf{SQL Server:}

\begin{minted}[fontsize=\small,frame=single,breaklines]{sql}
-- Phân trang bằng OFFSET
SELECT StoryId, Title, ChapterCount, LastChapterAt
FROM Story
WHERE PublishStatus = 'PUBLISHED' AND IsDeleted = 0 AND LastChapterAt IS NOT NULL
ORDER BY LastChapterAt DESC, StoryId DESC
OFFSET @offset ROWS FETCH NEXT 20 ROWS ONLY;

-- Phân trang bằng keyset: tiếp tục sau dòng cuối 
SELECT TOP (20) StoryId, Title, ChapterCount, LastChapterAt
FROM {table}
WHERE PublishStatus = 'PUBLISHED' AND IsDeleted = 0 AND LastChapterAt IS NOT NULL
  AND (LastChapterAt < ? OR (LastChapterAt = ? AND StoryId < ?))
ORDER BY LastChapterAt DESC, StoryId DESC

-- Chỉ mục có đúng thứ tự sắp xếp của truy vấn
CREATE NONCLUSTERED INDEX IX_Public_Last
    ON Story (PublishStatus, IsDeleted, LastChapterAt DESC, StoryId DESC)
    INCLUDE (Title, ChapterCount);
\end{minted}

\textbf{Redis:} một Sorted Set chứa các truyện công khai với điểm là thời điểm cập nhật; trang được lấy theo thứ hạng hoặc theo điểm. Biến thể RediSearch dùng \texttt{SORTBY} với \texttt{LIMIT}.

\begin{minted}[fontsize=\small,frame=single,breaklines]{text}
ZRANGE public:by_last <offset> <offset+19> REV                        -- theo thứ hạng
ZRANGE public:by_last <lastTime> -inf BYSCORE REV LIMIT 0 20         -- keyset
FT.SEARCH idx:story "@PublishStatus:{PUBLISHED} @IsDeleted:{0}"
    SORTBY LastChapterAt DESC LIMIT <offset> 20
\end{minted}

\begin{table}[H]
\centering\small
\begin{tabular}{|l|l|r|r|}
\hline
\textbf{DBMS} & \textbf{Biến thể} & \textbf{Trang đầu (ms)} & \textbf{Trang giữa (ms)} \\ \hline
SQL Server & Không chỉ mục (phải sắp xếp) & 2,995 & 2,991 \\ \hline
SQL Server & Chỉ mục + \texttt{OFFSET} & 0,050 & 0,330 \\ \hline
SQL Server & Chỉ mục + keyset & -- & 0,191 \\ \hline
Redis & Sorted Set theo thứ hạng & 0,075 & 0,064 \\ \hline
Redis & Sorted Set keyset & -- & 0,090 \\ \hline
Redis & RediSearch \texttt{SORTBY} & 0,174 & 0,339 \\ \hline
\end{tabular}
\caption{Kết quả kịch bản 3 (1\,000 truyện)}
\label{tab:idx-i3}
\end{table}

\textbf{Nhận xét.} Từ bảng kết quả, có thể nhận thấy khi chỉ mục có sẵn đúng thứ tự cần sắp xếp, SQL Server không còn phải sắp xếp toàn bộ dữ liệu và trả trang đầu nhanh hơn khoảng 60 lần. Với \texttt{OFFSET}, các dòng bị bỏ qua vẫn phải đọc nên trang càng sâu càng chậm; keyset tránh được điều này. Sorted Set của Redis cho thời gian gần như không đổi ở mọi độ sâu, trong khi \texttt{SORTBY} của RediSearch chậm dần giống \texttt{OFFSET}.

% -------------------------------------------------------------------
\paragraph{Kịch bản 4 -- Chỉ mục có lọc (filtered index).}

\textbf{Mục đích.} Trên thực tế, các truy vấn công khai chỉ đọc truyện đang công khai, nên không cần đưa truyện nháp, truyện bị ẩn hay đã xoá vào chỉ mục. Kịch bản so sánh chỉ mục trên toàn bộ dữ liệu với chỉ mục chỉ chứa phần dữ liệu cần dùng, về dung lượng và tốc độ lấy trang đầu tiên.

\textbf{SQL Server:}

\begin{minted}[fontsize=\small,frame=single,breaklines]{sql}
-- Chỉ mục đầy đủ: mọi truyện
CREATE NONCLUSTERED INDEX IX_Last_All
    ON Story (PublishStatus, IsDeleted, LastChapterAt DESC) INCLUDE (Title, ChapterCount);
-- Filtered index: chỉ truyện công khai
CREATE NONCLUSTERED INDEX IX_Last_Public
    ON Story (LastChapterAt DESC) INCLUDE (Title, ChapterCount)
    WHERE PublishStatus = 'PUBLISHED' AND IsDeleted = 0;
\end{minted}

\textbf{Redis:} Sorted Set chứa mọi truyện (phải lọc lại trạng thái khi đọc), Sorted Set chỉ chứa truyện công khai, và chỉ mục RediSearch khai báo với \texttt{FILTER}.

\begin{minted}[fontsize=\small,frame=single,breaklines]{text}
FT.CREATE idx:public ON HASH PREFIX 1 story:
    FILTER "@PublishStatus=='PUBLISHED' && @IsDeleted=='0'"
    SCHEMA LastChapterAt NUMERIC SORTABLE
\end{minted}

\begin{table}[H]
\centering\small
\begin{tabular}{|l|l|r|r|r|}
\hline
\textbf{DBMS} & \textbf{Biến thể} & \textbf{Số mục} & \textbf{Dung lượng (KB)} & \textbf{Thời gian (ms)} \\ \hline
SQL Server & Chỉ mục đầy đủ & 1\,000 & 104 & 0,047 \\ \hline
SQL Server & Filtered index & 856 & 80 & 0,036 \\ \hline
Redis & Sorted Set mọi truyện & 1\,000 & 93 & 0,113 \\ \hline
Redis & Sorted Set truyện công khai & 856 & 80 & 0,071 \\ \hline
Redis & RediSearch \texttt{FILTER} & 856 & 133 & 0,156 \\ \hline
\end{tabular}
\caption{Kết quả kịch bản 4 (1\,000 truyện)}
\label{tab:idx-i4}
\end{table}

\textbf{Nhận xét.} Nhìn vào kết quả, có thể thấy việc chỉ đánh chỉ mục phần dữ liệu cần truy vấn giúp chỉ mục nhỏ hơn khoảng 20\% và truy vấn nhanh hơn ở cả hai DBMS. SQL Server hỗ trợ trực tiếp bằng mệnh đề \texttt{WHERE} khi tạo chỉ mục; với Sorted Set, ứng dụng phải tự quyết định truyện nào được thêm hoặc xoá khỏi tập. Chỉ mục RediSearch tốn bộ nhớ nhiều hơn Sorted Set tự xây.

% -------------------------------------------------------------------
\paragraph{Kịch bản 5 -- Tìm kiếm toàn văn.}

\textbf{Mục đích.} Kịch bản này nhằm kiểm tra chức năng tìm truyện theo từ khoá trong tên và phần giới thiệu. Người dùng có thể gõ có dấu hoặc không dấu, nên ngoài tốc độ, kịch bản đo thêm \emph{recall}, tức tỉ lệ số truyện tìm được so với số truyện thực sự chứa từ khoá (261 truyện ở mức 1\,000 dòng).

\textbf{SQL Server:} so sánh tìm chuỗi con bằng \texttt{LIKE} (phải đọc mọi dòng) với full-text index dùng \texttt{CONTAINSTABLE}. Full-text index được tạo với bộ tách từ tiếng Việt trên catalog không phân biệt dấu.

\begin{minted}[fontsize=\small,frame=single,breaklines]{sql}
SELECT StoryId, Title FROM Story
WHERE Title LIKE N'%' + @word + N'%' OR Description LIKE N'%' + @word + N'%';

CREATE FULLTEXT INDEX ON Story (Title LANGUAGE 1066, Description LANGUAGE 1066)
    KEY INDEX PK_Story ON ftc_Novelle;
SELECT s.StoryId, s.Title
FROM Story s JOIN CONTAINSTABLE(Story, (Title, Description), @word) ft ON ft.[KEY] = s.StoryId;
\end{minted}

\textbf{Redis:} so sánh \texttt{SCAN} với chỉ mục RediSearch trường \texttt{TEXT}, và một biến thể trong đó ứng dụng lưu thêm trường \texttt{SearchText} là tên và giới thiệu đã bỏ dấu.

\begin{minted}[fontsize=\small,frame=single,breaklines]{text}
FT.CREATE idx:text ON HASH PREFIX 1 story: SCHEMA Title TEXT NOSTEM Description TEXT NOSTEM
FT.SEARCH idx:text "@Title|Description:(<từ khoá>)"
FT.CREATE idx:ascii ON HASH PREFIX 1 story: SCHEMA SearchText TEXT NOSTEM
FT.SEARCH idx:ascii "@SearchText:(<từ khoá đã bỏ dấu>)"
\end{minted}

\begin{table}[H]
\centering\small
\begin{tabular}{|l|l|r|r|r|}
\hline
\textbf{DBMS} & \textbf{Biến thể} & \textbf{Thời gian (ms)} & \textbf{Recall gõ có dấu} & \textbf{Recall gõ không dấu} \\ \hline
SQL Server & \texttt{LIKE} & 30,23 & 100\% & 0\% \\ \hline
SQL Server & Full-text \texttt{CONTAINSTABLE} & 3,33 & 100\% & 100\% \\ \hline
Redis & \texttt{SCAN} & 6,57 & 39\% & 0\% \\ \hline
Redis & RediSearch \texttt{TEXT} & 0,28 & 39\% & 0\% \\ \hline
Redis & RediSearch trường bỏ dấu & 0,41 & 100\% & 100\% \\ \hline
\end{tabular}
\caption{Kết quả kịch bản 5 (1\,000 truyện)}
\label{tab:idx-i5}
\end{table}

\textbf{Nhận xét.} Qua kết quả trên, có thể thấy full-text index của SQL Server nhanh hơn \texttt{LIKE} khoảng 9 lần và tìm đủ cả khi người dùng gõ không dấu, nhờ bộ tách từ tiếng Việt. RediSearch nhanh hơn về thời gian nhưng không hiểu dấu tiếng Việt, nên chỉ tìm được 39\% số truyện; muốn đạt 100\%, ứng dụng phải tự chuẩn hoá và lưu thêm một bản văn bản đã bỏ dấu.

% -------------------------------------------------------------------
\paragraph{Kịch bản 6 -- Chi phí ghi khi có chỉ mục.}

\textbf{Mục đích.} Như đã đề cập ở phần lý thuyết, chỉ mục làm đọc nhanh hơn nhưng mỗi lần ghi phải cập nhật thêm chỉ mục. Kịch bản đo thời gian thêm một truyện và cập nhật điểm đánh giá của một truyện khi bảng có 0 và 3 chỉ mục phụ (trên \texttt{RatingAvg}, \texttt{AuthorId} và \texttt{Slug} duy nhất).

\textbf{SQL Server:}

\begin{minted}[fontsize=\small,frame=single,breaklines]{sql}
INSERT INTO Story (StoryId, AuthorId, Title, Slug, ..., RatingAvg, ...) SELECT ...;
UPDATE Story SET RatingAvg = RatingAvg + 0.01 WHERE StoryId = @id;

CREATE NONCLUSTERED INDEX IX_Rating ON Story (RatingAvg);
CREATE NONCLUSTERED INDEX IX_Author ON Story (AuthorId);
CREATE UNIQUE NONCLUSTERED INDEX UX_Slug ON Story (Slug);
\end{minted}

\textbf{Redis:} với chỉ mục tự xây, ứng dụng phải ghi Hash và cập nhật cả ba cấu trúc chỉ mục trong cùng một giao dịch \texttt{MULTI}/\texttt{EXEC}; với RediSearch, chỉ cần ghi Hash và module tự cập nhật chỉ mục.

\begin{minted}[fontsize=\small,frame=single,breaklines]{text}
MULTI
HSET story:<id> ... RatingAvg <r>
ZADD by_rating <r> <id>
SADD author:<authorId> <id>
SET  slug:<slug> <id>
EXEC
\end{minted}

\begin{table}[H]
\centering\small
\begin{tabular}{|l|l|r|r|}
\hline
\textbf{DBMS} & \textbf{Biến thể} & \textbf{INSERT (ms)} & \textbf{UPDATE (ms)} \\ \hline
SQL Server & 0 chỉ mục phụ (đọc 6 / 2 trang) & 1,290 & 1,216 \\ \hline
SQL Server & 3 chỉ mục phụ (đọc 15 / 6 trang) & 1,507 & 1,313 \\ \hline
Redis & Không chỉ mục & 0,009 & 0,005 \\ \hline
Redis & 3 chỉ mục tự xây (\texttt{MULTI}) & 0,019 & 0,013 \\ \hline
Redis & 3 trường RediSearch & 0,037 & 0,021 \\ \hline
\end{tabular}
\caption{Kết quả kịch bản 6 (1\,000 truyện)}
\label{tab:idx-i7}
\end{table}

\textbf{Nhận xét.} Kết quả cho thấy với 3 chỉ mục phụ, lệnh \texttt{INSERT} của SQL Server phải đọc và ghi nhiều trang gấp 2,5 lần và chậm hơn khoảng 17\%. Trên Redis, chi phí ghi tăng 2--4 lần, nhưng giá trị tuyệt đối vẫn ở mức micro giây; đổi lại, với chỉ mục tự xây, việc giữ chỉ mục khớp với dữ liệu hoàn toàn do ứng dụng đảm nhận.

% -------------------------------------------------------------------
\subsubsection{So sánh và đánh giá}
\label{subsubsec:indexing-compare}

Kết quả của các kịch bản trên được tổng hợp lại theo từng tiêu chí trong Bảng~\ref{tab:idx-compare}, từ đó nhóm rút ra kết luận về cách sử dụng chỉ mục cho ứng dụng Novelle.

\begin{table}[H]
\centering\small
\renewcommand{\arraystretch}{1.25}
\begin{tabularx}{\textwidth}{|>{\raggedright\arraybackslash}p{3.2cm}|>{\raggedright\arraybackslash}X|>{\raggedright\arraybackslash}X|}
\hline
\textbf{Tiêu chí} & \textbf{MS SQL Server} & \textbf{Redis} \\ \hline
Tra cứu theo khoá / một cột & Nhanh khi có nonclustered index (khoảng 0,05~ms) & Nhanh nhất qua tên key hoặc index tự xây (khoảng 0,004~ms) \\ \hline
Truy vấn nhiều điều kiện & Bộ tối ưu tự chọn; cần covering index để thực sự được dùng & Phải tự thiết kế Sorted Set kết hợp hoặc dùng RediSearch \\ \hline
Sắp xếp và phân trang & Cần index đúng thứ tự; \texttt{OFFSET} chậm dần theo độ sâu & Sorted Set ổn định ở mọi độ sâu \\ \hline
Tìm toàn văn tiếng Việt & Full-text hỗ trợ tách từ và bỏ dấu, tìm đúng 100\% & Nhanh nhưng phải tự chuẩn hoá dấu \\ \hline
Chi phí duy trì chỉ mục & Tự động, nhất quán trong giao dịch & Tự xây thì ứng dụng phải tự đồng bộ; RediSearch tự động nhưng tốn bộ nhớ \\ \hline
Thời gian tuyệt đối & Mức trăm micro giây đến mili giây & Mức micro giây, nhanh hơn 3--240 lần \\ \hline
\end{tabularx}
\caption{So sánh cơ chế chỉ mục của SQL Server và Redis}
\label{tab:idx-compare}
\end{table}

\paragraph{Kết luận cho hệ thống Novelle.}
\begin{itemize}
    \item \textbf{SQL Server là nơi đặt các chỉ mục nghiệp vụ chính:} bộ lọc nhiều điều kiện, chỉ mục duy nhất và filtered index để bảo đảm ràng buộc, full-text index cho tìm kiếm truyện tiếng Việt. Chỉ mục được cập nhật tự động trong cùng giao dịch nên dữ liệu luôn nhất quán.
    \item \textbf{Redis dùng cho truy cập theo khoá với tần suất cao:} bộ nhớ đệm chi tiết truyện, chương, trang chủ và các bộ đếm, nơi tên key đã đóng vai trò chỉ mục và cho thời gian truy cập thấp nhất.
    \item Không xây dựng chỉ mục phụ trên Redis cho dữ liệu nghiệp vụ, vì ứng dụng phải tự duy trì tính nhất quán với SQL Server và RediSearch không xử lý tốt tiếng Việt có dấu.
\end{itemize}

\subsection{Xử lý truy vấn (Query Processing)}
\label{subsec:query-processing}

\subsubsection{Cơ sở lý thuyết}
% TODO

\subsubsection{Cơ chế xử lý truy vấn trên MS SQL Server}
% TODO

\subsubsection{Cơ chế xử lý lệnh trên Redis}
% TODO

\subsubsection{Thiết kế kịch bản thực nghiệm}
% TODO

\subsection{Xử lý giao tác (Transaction)}
\label{subsec:transaction}

\subsubsection{Cơ sở lý thuyết}
% TODO

\subsubsection{Cơ chế xử lý giao tác trên MS SQL Server}
% TODO

\subsubsection{Cơ chế xử lý giao tác trên Redis}
% TODO

\subsubsection{Thiết kế kịch bản và thực hiện thực nghiệm}
% TODO

\subsubsection{So sánh và đánh giá}
% TODO

\subsection{Sao lưu và phục hồi dữ liệu (Data Backup and Recovery)}
\label{subsec:backup-recovery}

\subsubsection{Cơ sở lý thuyết}
% TODO

\subsubsection{Cơ chế sao lưu và phục hồi dữ liệu trên MS SQL Server}
% TODO

\subsubsection{Cơ chế sao lưu và phục hồi dữ liệu trên Redis}
% TODO

\subsubsection{Thiết kế kịch bản và thực hiện thực nghiệm}
% TODO

\subsubsection{So sánh và đánh giá}
% TODO

\subsection{Tổng hợp và so sánh kết quả nghiên cứu}
\label{subsec:topics-summary}
% TODO


    \section{Thiết kế kiến trúc và triển khai ứng dụng}
\label{sec:architecture-implementation}

\subsection{Tổng quan công nghệ sử dụng}
% TODO

\subsection{Thiết kế kiến trúc hệ thống}
% TODO

\subsection{Thiết kế cơ chế tích hợp hai DBMS}
% TODO

\subsection{Hiện thực ứng dụng}

\subsubsection{Cấu trúc mã nguồn}
% TODO

\subsubsection{Hiện thực Backend}
% TODO

\subsubsection{Tích hợp MS SQL Server}
% TODO

\subsubsection{Tích hợp Redis}
% TODO

\subsubsection{Hiện thực Frontend}
% TODO

\subsubsection{Triển khai ứng dụng}
% TODO

\subsection{Minh họa các chức năng chính}
% TODO: chèn ảnh chụp màn hình vào image/chapter5/

    \section{Kết luận và đánh giá}
\label{sec:conclusion}

Chương cuối cùng tổng kết những gì nhóm đã thực hiện trong bài tập lớn, đặt lại hai hệ quản trị cạnh nhau dựa trên kết quả thực nghiệm của Chương~\ref{sec:dbms-topics}, đánh giá mức độ hoàn thành so với mục tiêu đặt ra ở Chương~\ref{sec:introduction}, và chỉ ra những hạn chế còn tồn tại của đề tài.

\subsection{Tổng kết kết quả thực hiện}
\label{subsec:conclusion-summary}

Bài tập lớn được thực hiện theo hai hướng gắn liền với nhau: nghiên cứu cơ chế bên trong của hai hệ quản trị, và vận dụng kết quả nghiên cứu vào việc xây dựng ứng dụng Novelle. Những kết quả chính mà nhóm đạt được có thể tóm tắt như sau.

\paragraph{Về phía nghiên cứu.} Nhóm đã tìm hiểu tổng quan về Microsoft SQL Server và Redis (Chương~\ref{sec:dbms-overview}), sau đó đi sâu vào năm chủ đề kỹ thuật: quản lý và lưu trữ dữ liệu, cấu trúc chỉ mục, xử lý truy vấn, xử lý giao tác, sao lưu và phục hồi dữ liệu (Chương~\ref{sec:dbms-topics}). Với mỗi chủ đề, nhóm trình bày cơ sở lý thuyết, phân tích cách mỗi hệ quản trị hiện thực vấn đề, rồi thiết kế các kịch bản thực nghiệm tương đương và chạy trên cùng một bộ dữ liệu của Novelle, gồm SQL Server 2019 trên Amazon RDS và Redis Stack 7.4. Các kịch bản không chỉ đo thời gian mà còn quan sát trực tiếp hành vi bên trong của hệ quản trị, chẳng hạn số trang LOB và độ phân mảnh chỉ mục, kế hoạch thực thi mà bộ tối ưu chọn ở từng kích thước dữ liệu, số deadlock khi tải tăng, số key bị trục xuất khi đầy bộ nhớ hay lượng dữ liệu được nạp lại sau sự cố. Nhờ vậy, mỗi nhận định so sánh trong báo cáo đều được giải thích bằng cơ chế và được kiểm chứng bằng số liệu.

\paragraph{Về phía ứng dụng.} Nhóm đã phân tích yêu cầu và xác định 13 nhóm yêu cầu chức năng cho bốn nhóm người dùng: khách, người đọc, tác giả và quản trị viên (Chương~\ref{sec:app-db-design}). Từ đó, nhóm xây dựng sơ đồ ERD, ánh xạ sang lược đồ quan hệ gồm 18 bảng trên SQL Server với đầy đủ khoá chính, khoá ngoại, ràng buộc duy nhất, chỉ mục, view, trigger và thủ tục lưu trữ; đồng thời thiết kế cách tổ chức dữ liệu trên Redis gồm quy ước đặt tên key, cấu trúc dữ liệu và thời gian sống cho sáu nhóm dữ liệu là bộ nhớ đệm, bộ đếm lượt xem, bản nháp, phiên đăng nhập, giới hạn tần suất và khoá phân tán. Ứng dụng được hiện thực theo kiến trúc nhiều tầng với frontend React và backend Spring Boot, tổ chức thành các module Nền tảng, Tài khoản, Danh mục \& truyện, Đọc truyện, Thư viện \& đánh giá, Yêu cầu tác giả, Studio sáng tác và Quản trị (Chương~\ref{sec:architecture-implementation}).

\paragraph{Mối liên hệ giữa nghiên cứu và ứng dụng.} Điểm nhóm đánh giá cao nhất ở đề tài là cách phân chia vai trò giữa hai hệ quản trị trong Novelle được rút ra từ kết quả thực nghiệm chứ không chọn theo cảm tính. Kết quả về lưu trữ cho thấy tự động lưu trên Redis nhanh hơn khoảng 250 lần, nên bản nháp được đặt trên Redis và chỉ ghi xuống SQL Server khi tác giả bấm Lưu. Kết quả về xử lý truy vấn và chỉ mục cho thấy SQL Server trả lời được mọi câu hỏi và tìm kiếm toàn văn tiếng Việt chính xác, nên tìm kiếm, lọc nhiều điều kiện và thống kê được giao cho SQL Server; còn các câu hỏi cố định được gọi rất nhiều như trang chủ, chi tiết truyện, nội dung chương được lưu đệm trên Redis. Kết quả về giao tác và sao lưu cho thấy SQL Server không mất giao tác đã commit và rollback đầy đủ khi có lỗi, nên mọi nghiệp vụ nhiều bước như duyệt tác giả, tạo chương hay đánh giá đều chạy trong giao tác trên SQL Server; còn bộ đếm lượt xem được tăng nguyên tử bằng script Lua trên Redis và đồng bộ dồn xuống SQL Server mỗi 60 giây, vì mất một phần nhỏ lượt xem khi có sự cố là chấp nhận được.

\subsection{Đánh giá và so sánh hai DBMS}
\label{subsec:conclusion-comparison}

Ở Chương~\ref{sec:dbms-overview}, nhóm đã đưa ra so sánh ban đầu rằng SQL Server ưu tiên tính toàn vẹn còn Redis ưu tiên tốc độ. Sau năm chủ đề nghiên cứu, nhận định này không chỉ được xác nhận mà còn được làm rõ hơn ở cấp độ cơ chế. Bảng~\ref{tab:conclusion-compare} tổng hợp kết quả chính của từng chủ đề.

\begin{table}[H]
\centering
\small
\renewcommand{\arraystretch}{1.3}
\begin{tabularx}{\textwidth}{|>{\raggedright\arraybackslash}p{2.6cm}|>{\raggedright\arraybackslash}X|>{\raggedright\arraybackslash}X|}
\hline
\textbf{Chủ đề} & \textbf{MS SQL Server} & \textbf{Redis} \\ \hline
Lưu trữ và quản lý dữ liệu & Dữ liệu nằm trên đĩa, RAM chỉ giữ bản sao của trang; khi vùng đệm quá nhỏ vẫn không mất dòng nào, nhưng đọc một chương mất khoảng 6,7~ms. & Dữ liệu nằm trên RAM, đọc một chương khoảng 0,002~ms, nhập hàng loạt nhanh hơn khoảng 18 lần; khi đầy bộ nhớ phải từ chối ghi hoặc trục xuất key. \\ \hline
Cấu trúc chỉ mục & Nhiều loại chỉ mục (clustered, nonclustered, covering, filtered, full-text), tự cập nhật trong giao tác; tìm toàn văn tiếng Việt đúng 100\%. & Chỉ có sẵn truy cập theo tên key; chỉ mục phụ phải tự xây và tự đồng bộ, hoặc dùng RediSearch nhưng không xử lý tốt tiếng Việt có dấu. \\ \hline
Xử lý truy vấn & Optimizer tự chọn và tự đổi kế hoạch theo dữ liệu (4/5 truy vấn đổi kế hoạch); truy vấn kết nối, truy vấn con, tổng hợp tăng tới khoảng 20--130~ms ở 10.000 truyện. & Không có kết nối, truy vấn con hay gom nhóm; ứng dụng chuẩn bị sẵn kết quả lúc ghi nên mọi truy vấn đều dưới 0,4~ms và không đổi theo dữ liệu. \\ \hline
Xử lý giao tác & Rollback đầy đủ, có ràng buộc làm lớp bảo vệ cuối cùng; nhưng thiết kế khoá sai gây deadlock (hơn 400 deadlock với trigger tính lại điểm). & Không có rollback và ràng buộc; đơn luồng nên không có deadlock, script Lua không xung đột; \texttt{WATCH} giảm mạnh khi tranh chấp cao. \\ \hline
Sao lưu và phục hồi & Không mất giao tác đã commit; giao tác chưa commit được rollback khi sự cố; khôi phục theo thời điểm nhờ log backup của RDS và temporal table. & Nạp lại đủ dữ liệu từ AOF sau khi khởi động lại; có thể mất khoảng một giây dữ liệu với \texttt{everysec}; AOF/RDB nằm cùng máy, không bảo vệ khi mất máy. \\ \hline
\end{tabularx}
\caption{Tổng hợp kết quả so sánh SQL Server và Redis qua năm chủ đề}
\label{tab:conclusion-compare}
\end{table}

Từ bảng tổng hợp, nhóm rút ra ba nhận xét chung.

\begin{itemize}
    \item \textbf{Hai hệ quản trị đặt công việc ở hai thời điểm khác nhau.} SQL Server \emph{làm việc lúc đọc}: người viết chỉ mô tả dữ liệu cần lấy, hệ quản trị tự tìm cách lấy và tự điều chỉnh khi dữ liệu thay đổi. Redis \emph{làm việc lúc ghi}: thứ tự, dữ liệu ghép và số liệu tổng hợp được chuẩn bị sẵn khi dữ liệu thay đổi, nên đọc luôn nhanh nhưng chỉ trả lời được những câu hỏi đã được thiết kế trước.
    \item \textbf{Trách nhiệm giữ dữ liệu đúng nằm ở hai nơi khác nhau.} Với SQL Server, phần lớn việc này do hệ quản trị đảm nhận thông qua ràng buộc, khoá, nhật ký giao tác và sao lưu. Với Redis, việc này chuyển sang ứng dụng: người lập trình phải tự duy trì chỉ mục phụ, tự gộp các lệnh vào script Lua và tự chọn chính sách trục xuất, cấu hình lưu bền phù hợp với từng loại dữ liệu.
    \item \textbf{Hai hệ quản trị bổ sung cho nhau chứ không thay thế nhau.} SQL Server phù hợp làm \textbf{nguồn dữ liệu gốc} cho dữ liệu nghiệp vụ nhiều quan hệ, cần truy vấn linh hoạt và cần đúng tuyệt đối. Redis phù hợp làm \textbf{lớp phục vụ nhanh} đứng phía trước, cho dữ liệu được truy cập thường xuyên, có dạng truy cập cố định và có thể tạo lại từ SQL Server hoặc chấp nhận mất một phần.
\end{itemize}

Bên cạnh đó, các thực nghiệm cũng cho thấy một hệ quản trị mạnh không tự động bảo đảm ứng dụng chạy tốt. Ở SQL Server, cùng là giao tác đúng nhưng trigger tính lại toàn bộ điểm đánh giá khiến thông lượng giảm xuống dưới 1 lượt mỗi giây ở 10 client, trong khi cách cộng dồn đạt khoảng 216 lượt mỗi giây ở 100 client mà không có deadlock nào. Ở Redis, cùng một giới hạn bộ nhớ nhưng đổi chính sách từ \texttt{noeviction} sang \texttt{allkeys-lru} làm tỷ lệ đọc trúng bộ nhớ đệm tăng từ khoảng 47\% lên khoảng 82\%, đổi lại 87,3\% lượt xem bị mất nếu bộ đếm đặt chung với bộ nhớ đệm. Nói cách khác, hiểu cơ chế bên trong của hệ quản trị là điều kiện cần để sử dụng nó đúng cách.

\subsection{Đánh giá mức độ hoàn thành đề tài}
\label{subsec:conclusion-completion}

Bảng~\ref{tab:conclusion-completion} đối chiếu các mục tiêu đặt ra ở Chương~\ref{sec:introduction} với kết quả thực tế mà nhóm đạt được.

\begin{table}[H]
\centering
\small
\renewcommand{\arraystretch}{1.3}
\begin{tabularx}{\textwidth}{|>{\raggedright\arraybackslash}p{4.2cm}|>{\raggedright\arraybackslash}X|>{\centering\arraybackslash}p{2.2cm}|}
\hline
\textbf{Mục tiêu} & \textbf{Kết quả đạt được} & \textbf{Mức độ} \\ \hline
Tìm hiểu tổng quan về SQL Server và Redis & Trình bày nguồn gốc, đặc điểm, lĩnh vực ứng dụng và so sánh tổng quan hai hệ quản trị (Chương~\ref{sec:dbms-overview}). & Hoàn thành \\ \hline
Nghiên cứu năm chủ đề kỹ thuật & Cả năm chủ đề đều có cơ sở lý thuyết, cơ chế trên từng hệ quản trị, kịch bản thực nghiệm tương đương, so sánh và kết luận cho Novelle (Chương~\ref{sec:dbms-topics}). & Hoàn thành \\ \hline
Kiểm chứng bằng thực nghiệm & Các kịch bản được chạy trên dữ liệu của Novelle ở nhiều mức kích thước; kết quả được ghi lại bằng số liệu, kế hoạch thực thi và biểu đồ. Một số phép đo còn giới hạn do môi trường (Mục~\ref{subsec:conclusion-limitations}). & Hoàn thành phần lớn \\ \hline
Phân tích yêu cầu và thiết kế cơ sở dữ liệu & 13 nhóm yêu cầu chức năng; ERD, lược đồ quan hệ 18 bảng trên SQL Server; thiết kế key, cấu trúc dữ liệu và thời gian sống trên Redis (Chương~\ref{sec:app-db-design}). & Hoàn thành \\ \hline
Xây dựng ứng dụng Novelle & Hiện thực đủ các module cho bốn nhóm người dùng, từ đọc truyện, thư viện, đánh giá, sáng tác chương đến duyệt tác giả và quản trị hệ thống (Chương~\ref{sec:architecture-implementation}). & Hoàn thành \\ \hline
Kết hợp hai hệ quản trị dựa trên kết quả nghiên cứu & SQL Server là nguồn dữ liệu gốc; Redis đảm nhận bộ nhớ đệm, bộ đếm lượt xem, bản nháp, phiên đăng nhập, giới hạn tần suất và khoá phân tán; mỗi lựa chọn đều được giải thích bằng kết quả thực nghiệm. & Hoàn thành \\ \hline
\end{tabularx}
\caption{Đánh giá mức độ hoàn thành các mục tiêu của đề tài}
\label{tab:conclusion-completion}
\end{table}

Nhìn chung, nhóm đã hoàn thành các mục tiêu chính của đề tài. Phần nghiên cứu không dừng lại ở việc liệt kê tính năng mà giải thích được \emph{vì sao} mỗi hệ quản trị hành xử như vậy và kiểm chứng bằng thực nghiệm; phần ứng dụng không chỉ dùng hai hệ quản trị cùng lúc mà phân chia vai trò cho từng loại dữ liệu dựa trên chính các kết quả đó. Thông qua đề tài, các thành viên cũng hiểu rõ hơn những khái niệm đã học trong môn Hệ quản trị cơ sở dữ liệu như tổ chức trang và vùng đệm, cây B\textsuperscript{+}, bộ tối ưu dựa trên chi phí, các thuật toán kết nối, khoá hai pha, deadlock và nhật ký giao tác, khi quan sát chúng hoạt động trên một hệ quản trị thực tế.

\subsection{Hạn chế của đề tài}
\label{subsec:conclusion-limitations}

Bên cạnh những kết quả đạt được, đề tài vẫn còn một số hạn chế, chủ yếu đến từ môi trường thực nghiệm và thời gian thực hiện.

\paragraph{Về môi trường và phạm vi thực nghiệm.}
\begin{itemize}
    \item \textbf{Phần cứng không tương đương:} SQL Server chạy trên Amazon RDS còn Redis chạy trên máy cá nhân, nên thời gian tuyệt đối và tỷ lệ tốc độ giữa hai hệ quản trị chỉ mang tính tham khảo; điều có giá trị hơn là xu hướng thay đổi theo kích thước dữ liệu và hành vi bên trong của từng hệ.
    \item \textbf{Quy mô dữ liệu còn nhỏ:} do giới hạn bộ nhớ của máy chủ RDS \texttt{db.t3.micro}, mức dữ liệu lớn nhất chỉ là 10.000 truyện. Dữ liệu ở mức 1.000 và 10.000 truyện được nhân bản từ 150 truyện thật nên phân bố giá trị đều hơn thực tế, có thể làm thống kê ước lượng của SQL Server dễ đúng hơn.
    \item \textbf{Một số phép đo chưa đầy đủ:} chi phí ghi của Redis khi phải cập nhật các cấu trúc dữ liệu chuẩn bị sẵn (Sorted Set, tổng rating, tổng theo ngày) chưa được đo; các số đo dưới 1~ms của Redis còn dao động giữa các lần chạy.
    \item \textbf{Chưa mô phỏng được sự cố thật:} nhóm không có quyền thực hiện native backup và phục hồi toàn instance trên RDS, nên RPO khi mất instance được suy ra từ lịch sử backup chứ không đo bằng một lần khôi phục thật. Phía Redis, lệnh \texttt{SHUTDOWN NOSAVE} vẫn xả AOF trước khi tắt nên chưa mô phỏng được tình huống mất điện.
\end{itemize}

\paragraph{Về ứng dụng.}
\begin{itemize}
    \item \textbf{Một số đề xuất từ nghiên cứu chưa được áp dụng hết:} trigger \texttt{trg\_Rating} vẫn tính lại toàn bộ điểm đánh giá thay vì cộng dồn, nên có nguy cơ deadlock khi nhiều người cùng đánh giá một truyện như kết quả ở Mục~\ref{subsubsec:tx-compare}.
    \item \textbf{Cấu hình Redis còn đơn giản:} hệ thống dùng một instance Redis cho mọi nhóm dữ liệu, trong khi thực nghiệm cho thấy bộ nhớ đệm và bộ đếm lượt xem nên được tách riêng hoặc dùng chính sách trục xuất chỉ áp dụng cho key có TTL. Ngoài ra, AOF và RDB nằm trên cùng máy với Redis nên chưa có bản sao lưu ở nơi khác.
    \item \textbf{Phạm vi chức năng có giới hạn:} ứng dụng tập trung vào các luồng cốt lõi của một nền tảng truyện; các tính năng như thanh toán, quảng cáo, thông báo thời gian thực hay gợi ý truyện theo sở thích chưa nằm trong phạm vi đề tài. Hệ thống cũng chưa được kiểm thử tải ở quy mô nhiều người dùng thật.
\end{itemize}

\paragraph{Hướng phát triển.} Từ những hạn chế trên, nhóm đề xuất một số hướng phát triển tiếp theo: chạy lại các thực nghiệm trên cùng một môi trường phần cứng với dữ liệu lớn hơn và có phân bố gần với thực tế hơn; bổ sung phép đo chi phí ghi của Redis và thử nghiệm khôi phục toàn instance; áp dụng cách cộng dồn cho \texttt{trg\_Rating}; tách Redis thành các instance riêng cho bộ nhớ đệm và dữ liệu cần giữ lâu hơn, kết hợp sao chép AOF/RDB ra kho lưu trữ độc lập; và mở rộng ứng dụng với các tính năng như thông báo hay gợi ý truyện, trong đó có thể tận dụng thêm các cấu trúc dữ liệu của Redis như Pub/Sub hoặc Stream.

    \cleardoublepage
    \nocite{elmasri2016fundamentals,silberschatz2019database,mssqldocs,redisdocs,carlson2013redis}
    \printbibliography[heading=bibintoc,title={Tài liệu tham khảo}]
    \appendix
    \cleardoublepage
    \phantomsection
    
    \section*{Phụ lục. Biên bản họp nhóm}
\addcontentsline{toc}{section}{Phụ lục. Biên bản họp nhóm}

% Ghi chú: 3 buổi họp chính (số 1, 2, 3) và 5 buổi họp phụ kiểm tra tiến độ,
% trải đều trong 7 tuần (28/09/2026 -- 16/11/2026).
% Có thể bổ sung giờ họp và nền tảng họp trực tuyến (Google Meet, MS Teams...) nếu cần.

% =====================================================================
\subsection*{Buổi họp số 1}

\begin{table}[H]
    \centering
    \renewcommand{\arraystretch}{1.4}
    \begin{tabular}{|>{\raggedright\arraybackslash}p{4cm}|>{\raggedright\arraybackslash}p{11.5cm}|}
        \hline
        \textbf{Thời gian} & 28/09/2026 \\
        \hline
        \textbf{Hình thức / Địa điểm} & Trực tuyến qua Google Meet \\
        \hline
        \textbf{Thành viên tham dự} & Toàn bộ thành viên nhóm \\
        \hline
        \textbf{Nội dung thảo luận} &
        \textbullet~Tìm hiểu yêu cầu của bài tập lớn và lựa chọn ứng dụng làm đề tài. \newline
        \textbullet~Lựa chọn hệ quản trị cơ sở dữ liệu. \\
        \hline
        \textbf{Kết luận và phân công} &
        \textbullet~Thống nhất đề tài: \textit{Ứng dụng đọc truyện trực tuyến}. \newline
        \textbullet~Mỗi thành viên tham khảo các ứng dụng hiện có để chuẩn bị cho buổi họp xác định chức năng. \\
        \hline
    \end{tabular}
\end{table}

% =====================================================================
\subsection*{Buổi họp số 2}

\begin{table}[H]
    \centering
    \renewcommand{\arraystretch}{1.4}
    \begin{tabular}{|>{\raggedright\arraybackslash}p{4cm}|>{\raggedright\arraybackslash}p{11.5cm}|}
        \hline
        \textbf{Thời gian} & 05/10/2026 \\
        \hline
        \textbf{Hình thức / Địa điểm} & Trực tuyến qua Google Meet \\
        \hline
        \textbf{Thành viên tham dự} & Toàn bộ thành viên nhóm \\
        \hline
        \textbf{Nội dung thảo luận} &
        \textbullet~Tổng hợp kết quả tham khảo các ứng dụng đọc truyện hiện có. \newline
        \textbullet~Xác định các chức năng cần thiết cho ứng dụng. \newline
        \textbullet~Xác định cấu trúc dữ liệu chính cần lưu trữ của hệ Thống. \\
        \hline
        \textbf{Kết luận và phân công} &
        \textbullet~Chốt danh sách các chức năng chính yếu của hệ thống lưu hành truyện trực tuyến giữa tác giả vả người đọc. \newline
        \textbullet~Phân công soạn đặc tả yêu cầu chức năng và phi chức năng. \newline
        \textbullet~Phân công thiết kế lược đồ cơ sở dữ liệu, từ điển dữ liệu và cấu trúc khóa. \\
        \hline
    \end{tabular}
\end{table}

% =====================================================================
\subsection*{Buổi họp số 3}

\begin{table}[H]
    \centering
    \renewcommand{\arraystretch}{1.4}
    \begin{tabular}{|>{\raggedright\arraybackslash}p{4cm}|>{\raggedright\arraybackslash}p{11.5cm}|}
        \hline
        \textbf{Thời gian} & Thứ Bảy, 10/10/2026 \\
        \hline
        \textbf{Hình thức / Địa điểm} & Trực tuyến qua Google Meet \\
        \hline
        \textbf{Thành viên tham dự} & Toàn bộ thành viên nhóm \\
        \hline
        \textbf{Nội dung thảo luận} &
        \textbullet~Rà soát bản đặc tả yêu cầu và thiết kế lược đồ cơ sở dữ liệu. \newline
        \textbullet~Thống nhất các ràng buộc, chỉ mục và cấu trúc khóa trên cơ sở dữ liệu. \newline
        \textbullet~Phân công nhiệm vụ thực hiện đề tài cho các thành viên. \\
        \hline
        \textbf{Kết luận và phân công} &
        \textbullet~Hoàn thiện từ điển dữ liệu và các nội dung bảng, chỉ mục cho cơ sở dữ liệu. \newline
        \textbullet~Thực hiện xây dựng ứng dụng đề tài đã lựa chọn. \\
        \hline
    \end{tabular}
\end{table}

% =====================================================================
\subsection*{Buổi họp số 4}

\begin{table}[H]
    \centering
    \renewcommand{\arraystretch}{1.4}
    \begin{tabular}{|>{\raggedright\arraybackslash}p{4cm}|>{\raggedright\arraybackslash}p{11.5cm}|}
        \hline
        \textbf{Thời gian} &  28/10/2026 \\
        \hline
        \textbf{Hình thức / Địa điểm} & Trực tuyến \\
        \hline
        \textbf{Thành viên tham dự} & Toàn bộ thành viên nhóm \\
        \hline
        \textbf{Nội dung thảo luận} &
        \textbullet~Xây dựng nội dung nghiên cứu và đối chiếu giữa 2 hệ quản trị cơ sở dữ liêu là MS SQL Server và Redis trong bối cảnh của đề tài. \newline
        \textbullet~Thống nhất các chủ đề nghiên cứu: data storage \& management, indexing, query processing, transaction, backup \& recovery. \newline
        \textbullet~Thống nhất công cụ và phương pháp kiểm thử trên tập dữ liệu thực tế của đề tài.\newline
        \textbullet~Kiểm tra tiến độ thực hiện đề tài.\\
        \hline
        \textbf{Kết luận và phân công} &
        \textbullet~Mỗi thành viên phụ trách một chủ đề nghiên cứu. \newline
        \textbullet~Tiếp tục thực hiện nội dung ứng dụng chính của đề tài.\\
        \hline
    \end{tabular}
\end{table}

% =====================================================================
\subsection*{Buổi họp số 5}

\begin{table}[H]
    \centering
    \renewcommand{\arraystretch}{1.4}
    \begin{tabular}{|>{\raggedright\arraybackslash}p{4cm}|>{\raggedright\arraybackslash}p{11.5cm}|}
        \hline
        \textbf{Thời gian} & 10/11/2026 \\
        \hline
        \textbf{Hình thức / Địa điểm} & Trực tuyến qua Google Meet\\
        \hline
        \textbf{Thành viên tham dự} & Toàn bộ thành viên nhóm \\
        \hline
        \textbf{Nội dung thảo luận} &
        \textbullet~Kiểm tra tiến độ và hoàn thiện ứng dụng đề tài. \newline
        \textbullet~Góp ý chỉnh sửa bổ sung cho ứng dụng. \\
        \hline
        \textbf{Kết luận và phân công} &
        \textbullet~Đẩy nhanh tiến độ và hoàn thiện giai đoạn cuối của ứng dụng. \\
        \hline
    \end{tabular}
\end{table}

% =====================================================================
\subsection*{Buổi họp 6}

\begin{table}[H]
    \centering
    \renewcommand{\arraystretch}{1.4}
    \begin{tabular}{|>{\raggedright\arraybackslash}p{4cm}|>{\raggedright\arraybackslash}p{11.5cm}|}
        \hline
        \textbf{Thời gian} & 17/11/2026 \\
        \hline
        \textbf{Hình thức / Địa điểm} & Trực tuyến qua Google Meet\\
        \hline
        \textbf{Thành viên tham dự} & Toàn bộ thành viên nhóm \\
        \hline
        \textbf{Nội dung thảo luận} &
        \textbullet~Kiểm tra việc hoàn tất kịch bản kiểm thử và tài liệu ghi chú của từng chủ đề. \newline
        \textbullet~Thống nhất khung sườn và nội dung thực hiện của bài báo cáo tổng hợp. \\
        \hline
        \textbf{Kết luận và phân công} &
        \textbullet~Tất cả thành viên cùng tham gia viết nội dung báo cáo cho ứng dụng đề tài.\newline
        \textbullet~Mỗi thành viên viết mục báo cáo cho chủ đề nghiên cứu của mình. \newline
        \textbullet~Bổ sung tài liệu tham khảo. \\
        \hline
    \end{tabular}
\end{table}


\end{document}
